\documentclass[11pt]{article}

% ============================================================
% Encoding and page layout
% ============================================================

\usepackage[T1]{fontenc}
\usepackage[utf8]{inputenc}
\usepackage[a4paper,margin=1.12in]{geometry}

% ============================================================
% Mathematics and theorem environments
% ============================================================

\usepackage{amsmath,amssymb,amsthm,mathtools}
\usepackage{microtype}
\usepackage{xcolor}
\usepackage{enumitem}
\usepackage{hyperref}
\usepackage{url}
\usepackage[nameinlink,capitalize]{cleveref}

\hypersetup{
    colorlinks=true,
    linkcolor=blue!55!black,
    citecolor=blue!55!black,
    urlcolor=blue!55!black
}

\newtheorem{theorem}{Theorem}[section]
\newtheorem{proposition}[theorem]{Proposition}
\newtheorem{lemma}[theorem]{Lemma}
\newtheorem{corollary}[theorem]{Corollary}
\newtheorem{definition}[theorem]{Definition}
\newtheorem{remark}[theorem]{Remark}
\newtheorem{assumption}[theorem]{Assumption}

% ============================================================
% Macros
% ============================================================

% ============================================================
% Basic macros for the inner-strip certification companion paper
% ============================================================

\newcommand{\ii}{\mathrm{i}}
\newcommand{\dd}{\mathrm{d}}
\newcommand{\eps}{\varepsilon}

\newcommand{\s}{s}
\newcommand{\sig}{\sigma}
\newcommand{\tauh}{\tau}

\newcommand{\NQ}{\mathcal N_Q}
\newcommand{\Beff}{\mathfrak B_{\rm eff}}
\newcommand{\Bvec}{\mathfrak B}
\newcommand{\Rtot}{\mathcal R_{\rm total}}

\newcommand{\Rband}{\mathcal R_{\rm band}}
\newcommand{\Rbd}{\mathcal R_{\rm bd}}
\newcommand{\Rfar}{\mathcal R_{\rm far}}
\newcommand{\Rend}{\mathcal R_{\rm end}}
\newcommand{\Rnonstat}{\mathcal R_{\rm nonstat}}
\newcommand{\Rfloor}{\mathcal R_{\rm floor}}
\newcommand{\Rcore}{\mathcal R_{\rm core}}

\newcommand{\Gbd}{G_{\rm bd}^{J}}
\newcommand{\Gfar}{G_{\rm far}^{J}}
\newcommand{\Gend}{G_{\rm end}^{J}}
\newcommand{\Gnonstat}{G_{\rm nonstat}^{J}}
\newcommand{\Gcore}{G_{\rm core}}

\newcommand{\PiQ}{\Pi_Q}
\newcommand{\Czero}{C_0}
\newcommand{\Cone}{C_1}
\newcommand{\Cinner}{C_{\rm inner}}
\newcommand{\gzero}{g_0}

\newcommand{\JSON}[1]{\texttt{#1}}
\newcommand{\code}[1]{\texttt{#1}}

\newcommand{\Qside}{\texttt{Q}}
\newcommand{\Jproj}{\texttt{J\_projected}}
\newcommand{\Zeroside}{\texttt{zero}}
\newcommand{\Mixedside}{\texttt{mixed}}

\newcommand{\CertFile}{\texttt{INNER\_STRIP\_CERTIFICATE\_SAFE.json}}

\newcommand{\RepoURL}{\url{https://github.com/vicblab/RH-inner-strip-certifier}}



\title{
A Certified Inner-Strip Two-Channel Quotient Obstruction
for the Riemann Zeta Function
}

\author{
V\'ictor Blanco Bataller
}

\date{\today}
\begin{document}
\maketitle
\begin{abstract}
This paper gives a machine-checkable certification of the inner-strip residual
estimates arising in the reflected two-channel quotient obstruction for the
Riemann zeta function.  The certified strip is the fixed central strip
\[
        |\sigma-\tfrac12|\le 0.01,
        \qquad s=\sigma+i\tau,
\]
with the critical line separated by a finite-width two-channel bracket
certificate.  The certified expansion has the form
\[
        \mathcal N_Q(s)
        =
        \frac{i}{\tau}\mathfrak B_{\rm eff}(\sigma,\tau)
        +
        \mathcal R_{\rm total}(s),
\]
where the total residual is decomposed into seven separately certified sectors,
\[
        \mathcal R_{\rm total}
        =
        \mathcal R_{\rm band}
        +
        \mathcal R_{\rm bd}
        +
        \mathcal R_{\rm far}
        +
        \mathcal R_{\rm end}
        +
        \mathcal R_{\rm nonstat}
        +
        \mathcal R_{\rm floor}
        +
        \mathcal R_{\rm core}.
\]
Each sector certificate contains both a value bound and a
\(\sigma\)-derivative bound, and the final strict merger also requires a
finite-width bracket certificate.  The final certified constants are
\[
        C_0
        =
        255675240189.76995047822220961036088344393198723387179813950427053690494808885261,
\]
\[
        C_1
        =
        291824585243.67624174026287399979132232317927055971912800732675440018563808683476,
\]
and therefore
\[
        C_{\rm inner}
        =
        \sqrt{C_0^2+C_1^2}
        =
        387983784453.345700134863573943992120362905925566406274034828268159
\]
Together with the certified bracket lower bound
\[
        g_0=0.2,
\]
the height inequality
\[
        \sqrt T>
        \frac{C_{\rm inner}}{g_0}\log T
\]
holds at \(T=10^{30}\), and in fact already holds for every
\(T\ge 10^{29}\).  Subject to the structural quotient identity and
two-channel obstruction theorem proved in the analytic companion paper, this
certifies the inner-strip exclusion of off-critical zeros in
\[
        0<|\Re s-\tfrac12|\le0.01
\]
for all
\[
        |\Im s|\ge10^{29}.
\]

Combined with the separately certified outer-strip theorem, this gives a
high-height exclusion throughout the critical strip away from the endpoint
neighborhoods of \(0\) and \(1\), at height \(|\Im s|\ge10^{30}\).
\end{abstract}

\section{Introduction}

This paper is the certification companion to the reflected two-channel quotient
obstruction for the Riemann zeta function.  The analytic part of the argument
reduces the exclusion of off-critical zeros in a fixed central strip to a
quantitative comparison between an explicit leading bracket vector and a
certified residual vector.  The purpose of the present paper is to record, in a
machine-checkable and reproducible form, the constants entering that comparison.

The quotient numerator used in the obstruction is denoted by
\[
        \mathcal N_Q(s),
        \qquad s=\sigma+i\tau.
\]
At a nontrivial zero of \(\zeta(s)\), the structural analytic setup gives
\[
        \mathcal N_Q(s)=0.
\]
The reflected two-channel expansion writes this numerator as
\[
        \mathcal N_Q(s)
        =
        \frac{i}{\tau}\mathfrak B_{\rm eff}(\sigma,\tau)
        +
        \mathcal R_{\rm total}(s).
\]
The term \(\mathfrak B_{\rm eff}\) is the explicit two-channel bracket.  The
residual term \(\mathcal R_{\rm total}\) contains all errors left after the
stationary, boundary, endpoint, band, nonstationary, floor, and algebraic core
pieces have been separated.  The certificate proves that this residual is
smaller than the bracket term at the certified height.

\begin{assumption}[Structural quotient interface]
The present paper assumes the structural quotient identity and the
two-channel obstruction interface established in the analytic companion paper.
In particular, at a nontrivial zero of \(\zeta(s)\) the quotient numerator
satisfies
\[
        \NQ(s)=0,
\]
and the certified expansion has the form
\[
        \NQ(s)
        =
        \frac{\ii}{\tauh}\Beff(\sig,\tauh)
        +
        \Rtot(s).
\]
The role of the present paper is to certify the bracket lower bound, the
sector residual bounds, the normalization ledger, and the final height
inequality for this interface.
\end{assumption}

The residual is decomposed as
\[
        \mathcal R_{\rm total}
        =
        \mathcal R_{\rm band}
        +
        \mathcal R_{\rm bd}
        +
        \mathcal R_{\rm far}
        +
        \mathcal R_{\rm end}
        +
        \mathcal R_{\rm nonstat}
        +
        \mathcal R_{\rm floor}
        +
        \mathcal R_{\rm core}.
\]
The seven labels correspond to:
\[
        {\rm band},\quad
        {\rm bd},\quad
        {\rm far},\quad
        {\rm end},\quad
        {\rm nonstat},\quad
        {\rm floor},\quad
        {\rm core}.
\]
The sector certificates are merged by a strict final verifier.  The merger
checks that all seven sector certificates are present, that the normalization
side of each sector is declared, that every sector contains both a value
constant and a \(\sigma\)-derivative constant, and that a finite-width bracket
certificate is present.  The final repository contains a one-command
reproduction script,
\[
        \texttt{./run\_all\_certification.sh},
\]
which generates the seven sector certificates, the bracket certificate, and the
final file
\[
        \texttt{INNER\_STRIP\_CERTIFICATE\_SAFE.json}.
\]
The final repository structure and tutorial-style reproduction instructions are
available in the public repository
\begin{center}
\RepoURL
\end{center}
The repository README records the intended strip, bracket parameters, target
height, expected merged constants, and one-command reproduction script
\[
        \code{./run\_all\_certification.sh}.
\]

The certification uses two residual channels.  The first is the even residual
channel, bounded by a constant \(C_0\).  The second is the
\(\sigma\)-derivative channel, bounded by a constant \(C_1\).  The final
merged certificate gives
\[
        C_0
        =
        255675240189.76995047822220961036088344393198723387179813950427053690494808885261,
\]
and
\[
        C_1
        =
        291824585243.67624174026287399979132232317927055971912800732675440018563808683476.
\]
Thus
\[
        C_{\rm inner}
        =
        \sqrt{C_0^2+C_1^2}
        =
        387983784453.34570013486357394399212036290592556640627403482826815918742177195616.
\]
The final JSON produced by the strict bracket-aware merger records this value
as
\[
        \texttt{C\_inner\_vector\_upper}.
\]
The strict merger used for the final theorem requires both the seven residual
sector files and the finite-width bracket certificate. 

The bracket certificate supplies the lower bound
\[
        g_0=0.2.
\]
In the certificate this is recorded as
\[
        \texttt{g0\_bracket\_vector\_lower = 0.2}.
\]
The finite-width certificate is based on the critical-line lower bound
\[
        g_{\rm crit}
        =
        \frac{\sqrt2}{48}(23\log2-2)
        =
        0.410781461979799978\ldots,
\]
together with a conservative compact Lipschitz audit on
\[
        0\le u\le0.01.
\]
The audit records
\[
        \|\mathfrak B(u,\tau)\|
        \ge
        g_{\rm crit}-20u.
\]
At \(u=0.01\), this lower bound remains strictly above \(0.2\).  The finite
bracket script and its certificate are included in the final reproduction
bundle. 

The certified obstruction condition is
\[
        \sqrt T>
        \frac{C_{\rm inner}}{g_0}\log T.
\]
Substituting the certified values gives
\[
        C_{\rm inner}
        =
        387983784453.34570013486357394399212036290592556640627403482826815918742177195616,
        \qquad
        g_0=0.2.
\]
Solving
\[
        \sqrt T=
        \frac{C_{\rm inner}}{g_0}\log T
\]
gives the numerical threshold
\[
        T_{\rm min}
        \approx
        1.5867647318922528\times10^{28}.
\]
Consequently, the clean power-of-ten height implied by the certificate is
\[
        T\ge10^{29}.
\]
The repository includes a height-threshold utility,
\[
        \code{tools/height\_threshold.py},
\]
which reproduces this computation.  The same repository also records the
one-command reproduction pipeline and the expected final constants
\[
        C_{\rm even}\approx2.5567524018976995\times10^{11},
        \qquad
        C_{\sigma}\approx2.9182458524367624\times10^{11},
\]
\[
        C_{\rm inner}\approx3.879837844533457\times10^{11}.
\]

We record the certified statement as follows.

\begin{theorem}[Certified inner-strip residual and height bound]
Let
\[
        s=\sigma+i\tau,
        \qquad
        \sigma\in[0.49,0.51],
        \qquad
        \tau\ge10.
\]
The certificate proves the residual decomposition
\[
        \mathcal R_{\rm total}
        =
        \mathcal R_{\rm band}
        +
        \mathcal R_{\rm bd}
        +
        \mathcal R_{\rm far}
        +
        \mathcal R_{\rm end}
        +
        \mathcal R_{\rm nonstat}
        +
        \mathcal R_{\rm floor}
        +
        \mathcal R_{\rm core},
\]
with bounds
\[
        |\mathcal R_{\rm even}(s)|
        \le
        C_0\frac{\log\tau}{\tau^{3/2}},
\]
and
\[
        |\partial_\sigma\mathcal R_{\rm even}(s)|
        \le
        C_1\frac{\log\tau}{\tau^{3/2}},
\]
where
\[
        C_0
        =
        255675240189.76995047822220961036088344393198723387179813950427053690494808885261,
\]
and
\[
        C_1
        =
        291824585243.67624174026287399979132232317927055971912800732675440018563808683476.
\]
Therefore
\[
        C_{\rm inner}
        =
        \sqrt{C_0^2+C_1^2}
        =
        387983784453.34570013486357394399212036290592556640627403482826815918742177195616.
\]
Together with the certified bracket lower bound
\[
        g_0=0.2,
\]
the inequality
\[
        \sqrt T>
        \frac{C_{\rm inner}}{g_0}\log T
\]
holds for \(T=10^{30}\), and more sharply for every \(T\ge10^{29}\).
\end{theorem}

\begin{corollary}[Certified central inner-strip exclusion]
Assume the structural quotient identity and two-channel obstruction theorem
proved in the analytic companion paper.  Then there are no off-critical
nontrivial zeros of \(\zeta(s)\) in the central inner strip
\[
        0<|\Re s-\tfrac12|\le0.01
\]
for
\[
        |\Im s|\ge10^{29}.
\]
In particular, the certified exclusion holds at the shared target height
\[
        |\Im s|\ge10^{30}.
\]
\end{corollary}

The remainder of the paper explains how the certificate is assembled.  The next
section describes the JSON schema, the sector normalization ledger, and the
distinction between \(Q\)-side, \(J\)-side projected, mixed, and zero sectors.
Subsequent sections document the bracket certificate, the boundary sector, the
band and endpoint sectors, the far and global nonstationary sectors, and the
final merge.


% ============================================================
% Step 2: Certificate schema and normalization ledger
% ============================================================

\section{Certificate schema and normalization ledger}

The purpose of the certificate is not only to give numerical constants, but to
record a reproducible and auditable decomposition of the residual term
\[
        \Rtot
        =
        \Rband
        +
        \Rbd
        +
        \Rfar
        +
        \Rend
        +
        \Rnonstat
        +
        \Rfloor
        +
        \Rcore .
\]
For this reason the final certifier does not accept a single undifferentiated
constant.  Instead, it requires seven separate sector certificates and one
finite-width bracket certificate.  Each sector certificate records the sector
name, its normalization side, its value-channel constant, its
\(\sig\)-derivative-channel constant, a coverage hash, a normalization hash,
and a proof hash.  The strict final merger checks that all required sectors are
present and that the bracket certificate is present before producing the final
file
\[
        \CertFile .
\]

The final certificate is generated by the bracket-aware merger
\[
        \code{merge\_inner\_certificate\_with\_bracket\_safe.py}.
\]
This merger requires the seven residual-sector certificates
\[
        \code{band\_sector.safe.json},\quad
        \code{bd\_sector.safe.json},\quad
        \code{far\_sector.safe.json},\quad
        \code{end\_sector.safe.json},
\]
\[
        \code{nonstat\_sector.safe.json},\quad
        \code{floor\_sector.safe.json},\quad
        \code{core\_sector.safe.json},
\]
and the finite-width bracket certificate
\[
        \code{inner\_bracket\_vector.safe.json}.
\]
The merger refuses to produce a final theorem-level certificate if any of these
files is missing, if any residual sector does not have status
\[
        \JSON{proved},
\]
or if the bracket certificate does not have status
\[
        \JSON{proved}.
\]
This prevents the final theorem from being certified using only the limiting
critical-line bracket bound.

Each residual-sector certificate has the following logical schema:
\[
\begin{aligned}
        &\JSON{status},\\
        &\JSON{sector},\\
        &\JSON{side},\\
        &\JSON{C\_value\_upper\_safe},\\
        &\JSON{C\_sigma\_derivative\_upper\_safe},\\
        &\JSON{coverage\_hash},\\
        &\JSON{normalization\_hash},\\
        &\JSON{proof\_hash}.
\end{aligned}
\]
The field
\[
        \JSON{status}
\]
must equal
\[
        \JSON{proved}.
\]
The field
\[
        \JSON{sector}
\]
must be one of
\[
        \JSON{band},\quad
        \JSON{bd},\quad
        \JSON{far},\quad
        \JSON{end},\quad
        \JSON{nonstat},\quad
        \JSON{floor},\quad
        \JSON{core}.
\]
The field
\[
        \JSON{side}
\]
declares how the sector is normalized.  The allowed values are
\[
        \Qside,\qquad
        \Jproj,\qquad
        \Zeroside,\qquad
        \Mixedside .
\]

The meaning of these four normalization types is as follows.  A sector marked
\[
        \Qside
\]
already lives in quotient-numerator normalization.  Such a sector must not be
multiplied by the quotient projection again.  This is the case for the band
sector:
\[
        \Rband
        =
        C_B S_2^{\rm band}
        +
        \frac{\s-1}{\s}C_A S_1^{\rm band}
        -
        \frac{\ii}{2\tauh}
        \left(
            2^{-\sig}e^{-\ii L}
            -
            2^{\sig-1}e^{\ii L}
        \right).
\]
The band term is therefore declared as
\[
        \JSON{side = Q}.
\]

A sector marked
\[
        \Jproj
\]
is first naturally bounded on the \(J\)-side and is then projected into the
quotient numerator by
\[
        \PiQ(\s)
        =
        \frac{\ii\tauh}{(1-\s)(\s+1)}.
\]
For such a sector, if the \(J\)-side residual is \(G_\nu^{J}(\s)\), the
quotient-side residual is
\[
        \mathcal R_\nu(\s)
        =
        \PiQ(\s)G_\nu^{J}(\s).
\]
The derivative certificate must include the derivative of the projection:
\[
        \partial_\sig \mathcal R_\nu(\s)
        =
        \left(\partial_\sig \PiQ(\s)\right)G_\nu^{J}(\s)
        +
        \PiQ(\s)\partial_\sig G_\nu^{J}(\s),
\]
where
\[
        \partial_\sig\PiQ(\s)
        =
        \PiQ(\s)
        \left(
            \frac{1}{1-\s}
            -
            \frac{1}{\s+1}
        \right).
\]
The boundary, far, endpoint, and global nonstationary sectors are all declared
as
\[
        \JSON{side = J\_projected}.
\]

The floor sector is declared as
\[
        \JSON{side = zero}.
\]
It is identically zero under the half-open floor convention:
\[
        \Rfloor(\s)=0,
        \qquad
        \partial_\sig\Rfloor(\s)=0.
\]
Accordingly, its safe value and derivative constants are both zero.

The core sector is declared as
\[
        \JSON{side = mixed}.
\]
Its primary explicit contribution is the quotient-side coefficient mismatch
\[
        \Delta_\rho(\s,\tauh)
        =
        -\ii
        \left(\frac{2\pi}{\tauh}\right)^{1-2\sig}
        e^{2\ii\Theta(\tauh)}
        \frac{\s^2-\s+1}{\s(1-\s)(\s+1)}.
\]
The corresponding core component is
\[
        G_{\rho}^{\rm core}(\s)
        =
        C_B(\s,\tauh)\Delta_\rho(\s,\tauh)
        \sum_{m\le N(\tauh)} m^{-\s}.
\]
This term is already quotient-side and is not projected again.  If later
algebraic subcomponents are split off on the \(J\)-side, they must be declared
inside the core certificate and projected separately.  The certificate therefore
marks the whole core sector as mixed and records the projection convention in
the normalization hash.

Every raw sector constant is converted into a safety-normalized constant before
it is accepted by the final merger.  The safety rule is
\[
        C_{\rm safe}
        =
        \begin{cases}
        0, & C_{\rm raw}=0,\\
        C_{\rm raw}(1+10^{-12})+100, & C_{\rm raw}>0.
        \end{cases}
\]
The fields
\[
        \JSON{C\_value\_upper\_safe}
        \quad\text{and}\quad
        \JSON{C\_sigma\_derivative\_upper\_safe}
\]
are the only constants summed by the final merger.  This rule gives a uniform
rounding and transcription margin and prevents marginal certificates from
depending on exact decimal formatting.

The two hashes
\[
        \JSON{coverage\_hash}
        \quad\text{and}\quad
        \JSON{normalization\_hash}
\]
serve different purposes.  The coverage hash records what region or subregion
the sector certifies.  For example, the boundary sector coverage hash depends
on the compact-boundary hook, the midrange compact bridge, the high-\(\tauh\)
Taylor hook, the compact-complement tail hook, the compact-complement
nonstationary hook, and the stationary-overlap hook.  The normalization hash
records whether the sector is \(Q\)-side, \(J\)-side projected, zero, or mixed,
and therefore records whether \(\PiQ\) and \(\partial_\sig\PiQ\) are included.
Finally, the proof hash is a hash of the complete JSON object.  The final merge
records the proof hash of every sector and the proof hash of the bracket
certificate.

The boundary sector has its own internal ledger because it is assembled from
several hooks.  In the final run it is the sum of:
\[
        \text{finite compact boundary},
\]
\[
        \text{midrange compact boundary bridge},
\]
\[
        \text{high-\(\tauh\) compact Taylor boundary},
\]
\[
        \text{compact-complement tail},
\]
\[
        \text{compact-complement nonstationary transition},
\]
and
\[
        \text{compact-complement stationary overlap}.
\]
The boundary merger sums the safety-normalized constants from these hooks and
produces the final file
\[
        \code{bd\_sector.safe.json}.
\]
Both the value channel and the \(\sig\)-derivative channel are summed.  The
stationary and nonstationary transition hooks use complementary box counts:
boxes for which the phase derivative is separated from zero are assigned to the
nonstationary hook, and boxes for which the derivative interval intersects the
stationary threshold are assigned to the stationary-overlap hook.  This prevents
the transition region from being silently omitted.

The final strict merger then forms
\[
        \Czero
        =
        \sum_\nu C_{\nu,0}^{\rm safe},
        \qquad
        \Cone
        =
        \sum_\nu C_{\nu,1}^{\rm safe},
\]
where \(\nu\) runs over
\[
        {\rm band},\,
        {\rm bd},\,
        {\rm far},\,
        {\rm end},\,
        {\rm nonstat},\,
        {\rm floor},\,
        {\rm core}.
\]
It then computes
\[
        \Cinner
        =
        \sqrt{\Czero^2+\Cone^2}.
\]
For the final certificate produced in this run, the merged values are
\[
        \Czero
        =
        255675240189.76995047822220961036088344393198723387179813950427053690494808885261,
\]
\[
        \Cone
        =
        291824585243.67624174026287399979132232317927055971912800732675440018563808683476,
\]
and
\[
        \Cinner
        =
        387983784453.34570013486357394399212036290592556640627403482826815918742177195616.
\]
The merger then checks the height inequality
\[
        \sqrt T>
        \frac{\Cinner}{\gzero}\log T.
\]
The final theorem-level certificate is issued only if this inequality holds and
only if the bracket certificate proves
\[
        \gzero\ge0.2.
\]

Thus the certificate schema is deliberately redundant.  It records the sector
constant, the derivative constant, the sector side, the coverage declaration,
the normalization declaration, and the proof hash.  The resulting final JSON is
not merely a numerical value of \(\Cinner\); it is a reproducible dependency
ledger for the inner-strip obstruction.


% ============================================================
% Step 3: Finite-width bracket certificate
% ============================================================

\section{The finite-width bracket certificate}

The residual estimates certified in the previous section are useful only after
they are compared against a strictly positive leading bracket.  The analytic
obstruction separates the quotient numerator into the form
\[
        \NQ(s)
        =
        \frac{\ii}{\tauh}\Beff(\sig,\tauh)
        +
        \Rtot(s).
\]
At a zero of \(\zeta(s)\), one has \(\NQ(s)=0\).  Thus a contradiction follows
once the effective bracket is bounded below and the residual vector is bounded
above.  The present section records the finite-width lower bound used by the
strict final merger.

We write
\[
        \sig=\frac12+u,
        \qquad
        0\le u\le u_0,
        \qquad
        u_0=0.01.
\]
The phase
\[
        L=\tauh\log 2
\]
is treated modulo \(2\pi\), and the endpoint scalar
\[
        X=\Xi_{\rm end}(\tauh)
\]
is treated as a real parameter.  In the certificate the bracket is therefore
considered as a finite-dimensional real problem in
\[
        (u,L,X)\in[0,0.01]\times(\mathbb R/2\pi\mathbb Z)\times\mathbb R.
\]
The effective bracket appearing in the quotient expansion is
\[
\begin{aligned}
        \Beff(\sig,\tauh)
        =
        &\,
        2^\sig e^{\ii L}
        \left(
            \frac{12-\sig}{48}
            +
            \frac{\ii}{4}X
        \right)
        \\
        &-
        2^{-\sig}e^{-\ii L}
        \left(
            \frac{11+\sig}{24}
            -
            \frac{\ii}{2}X
        \right).
\end{aligned}
\]
After substituting \(\sig=\tfrac12+u\), this becomes
\[
\begin{aligned}
        \Beff(u,L,X)
        =
        &\,
        2^{1/2+u} e^{\ii L}
        \left(
            \frac{23/2-u}{48}
            +
            \frac{\ii}{4}X
        \right)
        \\
        &-
        2^{-1/2-u} e^{-\ii L}
        \left(
            \frac{23/2+u}{24}
            -
            \frac{\ii}{2}X
        \right).
\end{aligned}
\]
The two-channel obstruction uses the bracket vector obtained from the even
channel and the \(\sig\)-derivative channel.  We denote this real vector by
\[
        \Bvec(u,L,X).
\]
The quantity to be certified is the minimized bracket size
\[
        G(u)
        :=
        \inf_{L\in\mathbb R/2\pi\mathbb Z}
        \inf_{X\in\mathbb R}
        \|\Bvec(u,L,X)\|.
\]
The finite-width bracket certificate proves
\[
        G(u)\ge 0.2
\]
uniformly for
\[
        0\le u\le0.01.
\]

The minimization in \(X\) is not handled by sampling an unbounded real line.
For fixed \(u\) and \(L\), the squared bracket norm is a quadratic polynomial
in \(X\):
\[
        \|\Bvec(u,L,X)\|^2
        =
        A(u,L)X^2+B(u,L)X+C(u,L),
\]
with
\[
        A(u,L)>0.
\]
Thus the \(X\)-minimum is explicit:
\[
        X_*(u,L)
        =
        -\frac{B(u,L)}{2A(u,L)}.
\]
Substituting \(X_*(u,L)\) gives the reduced squared bracket
\[
        \Phi(u,L)
        :=
        C(u,L)
        -
        \frac{B(u,L)^2}{4A(u,L)}.
\]
Consequently
\[
        G(u)^2
        =
        \inf_{L\in\mathbb R/2\pi\mathbb Z}
        \Phi(u,L).
\]
This reduction is important: the actual finite-width certificate is a compact
interval problem in the two variables
\[
        u\in[0,0.01],
        \qquad
        L\in[0,2\pi].
\]
The real parameter \(X\) has already been minimized analytically.

On the critical line \(u=0\), the reduced bracket admits the explicit lower
bound
\[
        G(0)
        \ge
        g_{\rm crit},
\]
where
\[
        g_{\rm crit}
        =
        \frac{\sqrt2}{48}(23\log2-2).
\]
Numerically,
\[
        g_{\rm crit}
        =
        0.410781461979799978\ldots .
\]
This is well above the final target
\[
        g_0=0.2.
\]
However, the final theorem is not a statement only on the critical line.  It
requires a finite-width lower bound on the whole interval
\[
        0\le u\le0.01.
\]
For this reason the certificate also records a compact Lipschitz audit for the
reduced bracket.

The finite-width audit proves the conservative bound
\[
        |G(u)-G(0)|
        \le
        20u
\]
on
\[
        0\le u\le0.01.
\]
Equivalently,
\[
        G(u)
        \ge
        G(0)-20u.
\]
Combining this with the critical-line lower bound gives
\[
        G(u)
        \ge
        g_{\rm crit}-20u.
\]
At the endpoint \(u=0.01\), this yields
\[
        G(u)
        \ge
        0.410781461979799978\ldots - 0.2
        =
        0.210781461979799978\ldots .
\]
In particular,
\[
        G(u)>0.2
\]
throughout
\[
        0\le u\le0.01.
\]
The certificate records the public lower bound as
\[
        g_0=0.2.
\]

The Lipschitz constant \(20\) is deliberately conservative.  The relevant
derivatives are elementary combinations of the functions
\[
        2^{1/2+u},
        \qquad
        2^{-1/2-u},
\]
the linear coefficients
\[
        \frac{23/2-u}{48},
        \qquad
        \frac{23/2+u}{24},
\]
and the reduced quadratic expression
\[
        \Phi(u,L)
        =
        C(u,L)-\frac{B(u,L)^2}{4A(u,L)}.
\]
The audit is performed after the analytic minimization over \(X\).  Hence the
only compact domain to be checked is
\[
        [0,0.01]\times[0,2\pi].
\]
The finite-width bracket certificate records this as a compact Lipschitz audit.
After the analytic minimization in \(X\), the remaining domain is compact:
\[
        (u,L)\in[0,0.01]\times[0,2\pi].
\]
The certificate bounds the \(u\)-variation of the reduced minimized bracket by
the conservative constant
\[
        K_{\rm br}=20.
\]
Equivalently,
\[
        G(u)\ge G(0)-K_{\rm br}u.
\]
Since
\[
        G(0)\ge
        \frac{\sqrt2}{48}(23\log2-2)
        =
        0.410781461979799978\ldots,
\]
one obtains at \(u=0.01\)
\[
        G(u)\ge
        0.410781461979799978\ldots-0.2
        =
        0.210781461979799978\ldots
        >
        0.2.
\]
Thus the public bound
\[
        g_0=0.2
\]
has a positive numerical margin.

The bracket certificate is emitted as the file
\[
        \code{certs\_inner/inner\_bracket\_vector.safe.json}.
\]
Its essential fields are
\[
        \JSON{status = proved},
\]
\[
        \JSON{sector = inner\_bracket\_vector},
\]
\[
        \JSON{u0 = 0.01},
\]
and
\[
        \JSON{g0\_bracket\_vector\_lower = 0.2}.
\]
The certificate also records the critical-line lower bound, the Lipschitz
constant, the computed endpoint lower bound, and a proof hash.  In the final
run the bracket dependency was recorded by the strict merger as
\[
        \JSON{bracket\_certificate}.
\]
The strict merger refuses to produce a theorem-level certificate unless this
file is present and has status
\[
        \JSON{proved}.
\]
Thus the final JSON cannot be produced from the residual certificates alone.
It requires both the residual-sector ledger and the finite-width bracket
certificate.  The bracket script and the strict bracket-aware final merger are
part of the final certifier repository. 

\begin{lemma}[Bracket margin]
The finite-width bracket certificate has numerical margin
\[
        g_{\rm crit}-20(0.01)-0.2
        =
        0.010781461979799978\ldots
        >
        0.
\]
Consequently the certified lower bound \(g_0=0.2\) is not a marginal equality.
\end{lemma}

\begin{proof}
The certificate records
\[
        g_{\rm crit}
        =
        \frac{\sqrt2}{48}(23\log2-2)
        =
        0.410781461979799978\ldots.
\]
The compact Lipschitz audit gives loss at most
\[
        20u\le20(0.01)=0.2.
\]
Therefore
\[
        G(u)\ge0.410781461979799978\ldots-0.2
        =
        0.210781461979799978\ldots,
\]
which exceeds \(0.2\) by approximately \(1.0781\times10^{-2}\).
\end{proof}

We summarize the finite-width certificate as follows.

\begin{theorem}[Finite-width bracket certificate]
Let
\[
        0\le u\le0.01.
\]
Let
\[
        G(u)
        =
        \inf_{L\in\mathbb R/2\pi\mathbb Z}
        \inf_{X\in\mathbb R}
        \|\Bvec(u,L,X)\|
\]
be the minimized two-channel bracket norm.  The certificate proves
\[
        G(u)\ge0.2
\]
uniformly on the interval
\[
        0\le u\le0.01.
\]
More precisely,
\[
        G(0)
        \ge
        \frac{\sqrt2}{48}(23\log2-2),
\]
and the finite-width audit proves
\[
        |G(u)-G(0)|\le20u.
\]
Therefore
\[
        G(u)
        \ge
        \frac{\sqrt2}{48}(23\log2-2)-20u
        >
        0.2
\]
for all
\[
        0\le u\le0.01.
\]
\end{theorem}

\begin{proof}
For fixed \(u\) and \(L\), the squared two-channel bracket norm is a positive
quadratic polynomial in the endpoint parameter \(X\).  Therefore \(X\) is
minimized exactly by
\[
        X_*(u,L)
        =
        -\frac{B(u,L)}{2A(u,L)}.
\]
Substitution gives the reduced compact problem
\[
        G(u)^2
        =
        \inf_{L\in[0,2\pi]}
        \left(
            C(u,L)
            -
            \frac{B(u,L)^2}{4A(u,L)}
        \right).
\]
On \(u=0\), the reduced expression gives the explicit critical-line lower
bound
\[
        G(0)
        \ge
        \frac{\sqrt2}{48}(23\log2-2).
\]
The interval audit then bounds the variation of the minimized bracket on
\[
        [0,0.01]\times[0,2\pi]
\]
by
\[
        |G(u)-G(0)|\le20u.
\]
Hence
\[
        G(u)
        \ge
        \frac{\sqrt2}{48}(23\log2-2)-20u.
\]
At \(u=0.01\), the right-hand side is still larger than \(0.2\).  This proves
the claimed finite-width lower bound.
\end{proof}

% ============================================================
% Step 4: Boundary sector certificate
% ============================================================

\section{The boundary sector certificate}

The boundary sector is the largest certified sector in the inner-strip
certificate.  It accounts for the stationary transition occurring near the
endpoint
\[
        x=2
\]
in the two stationary four-family tails.  In the structural decomposition this
sector is denoted by
\[
        \Rbd .
\]
Unlike the band sector, which is already naturally in quotient normalization,
the boundary sector is first controlled on the \(J\)-side and is then projected
into quotient normalization.  Thus the boundary certificate is declared with
\[
        \JSON{side = J\_projected}.
\]
If the \(J\)-side boundary residual is denoted by
\[
        \Gbd(\s),
\]
then the quotient-side residual is
\[
        \Rbd(\s)
        =
        \PiQ(\s)\Gbd(\s),
\]
where
\[
        \PiQ(\s)
        =
        \frac{\ii\tauh}{(1-\s)(\s+1)}.
\]
The derivative certificate therefore includes
\[
        \partial_\sig \Rbd(\s)
        =
        \left(\partial_\sig\PiQ(\s)\right)\Gbd(\s)
        +
        \PiQ(\s)\partial_\sig \Gbd(\s),
\]
with
\[
        \partial_\sig\PiQ(\s)
        =
        \PiQ(\s)
        \left(
            \frac{1}{1-\s}
            -
            \frac{1}{\s+1}
        \right).
\]
This projection derivative is included in the boundary hook constants.

The boundary analysis uses the scaled endpoint variables
\[
        w=\tauh^{-1/2},
        \qquad
        \alpha=\frac{\tauh}{4\pi},
        \qquad
        r_m=w(m-\alpha),
\]
and
\[
        x=2+wy.
\]
The endpoint layer is therefore studied in the compactified \(y\)-coordinate.
The standard compact cutoff is
\[
        \chi_Y(y),
        \qquad
        Y=128,
\]
where
\[
        \chi_Y(y)=1
        \quad
        (0\le y\le Y),
\]
\[
        \chi_Y(y)=1-3t^2+2t^3,
        \qquad
        t=\frac{y-Y}{Y},
        \quad
        (Y\le y\le 2Y),
\]
and
\[
        \chi_Y(y)=0
        \quad
        (y\ge2Y).
\]
The compact boundary integrals are therefore supported on
\[
        0\le y\le2Y.
\]

The two boundary phases are
\[
        \Phi_-(r,y,w)
        =
        2\pi r y
        +
        \frac{y}{2w}
        -
        \frac{1}{w^2}
        \log\left(1+\frac{wy}{2}\right),
\]
and
\[
        \Phi_+(r,y,w)
        =
        -\Phi_-(r,y,w).
\]
At \(w=0\), these have the limiting forms
\[
        \Phi_-(r,y,0)
        =
        2\pi r y+\frac{y^2}{8},
\]
and
\[
        \Phi_+(r,y,0)
        =
        -2\pi r y-\frac{y^2}{8}.
\]
The amplitudes are
\[
        A_-(\sig,y,w)
        =
        (2+wy)^{-\sig-2},
\]
and
\[
        A_+(\sig,y,w)
        =
        (2+wy)^{\sig-3}.
\]
The compact finite-\(w\) difference integrals are
\[
\begin{aligned}
        D_-^\chi(\sig,r,w)
        =
        \int_0^{2Y}
        \chi_Y(y)
        \bigg[
            &(2+wy)^{-\sig-2}
            e^{\ii\Phi_-(r,y,w)}
            \\
            &-
            2^{-\sig-2}
            e^{\ii\Phi_-(r,y,0)}
        \bigg]
        \,\dd y,
\end{aligned}
\]
and
\[
\begin{aligned}
        D_+^\chi(\sig,r,w)
        =
        \int_0^{2Y}
        \chi_Y(y)
        \bigg[
            &(2+wy)^{\sig-3}
            e^{\ii\Phi_+(r,y,w)}
            \\
            &-
            2^{\sig-3}
            e^{\ii\Phi_+(r,y,0)}
        \bigg]
        \,\dd y.
\end{aligned}
\]
The \(\sig\)-derivatives are obtained by differentiating only the amplitudes,
because the phases are independent of \(\sig\).  Therefore
\[
\begin{aligned}
        \partial_\sig D_-^\chi(\sig,r,w)
        =
        \int_0^{2Y}
        \chi_Y(y)
        \bigg[
            &-\log(2+wy)(2+wy)^{-\sig-2}
            e^{\ii\Phi_-(r,y,w)}
            \\
            &+
            \log2\,2^{-\sig-2}
            e^{\ii\Phi_-(r,y,0)}
        \bigg]
        \,\dd y,
\end{aligned}
\]
and
\[
\begin{aligned}
        \partial_\sig D_+^\chi(\sig,r,w)
        =
        \int_0^{2Y}
        \chi_Y(y)
        \bigg[
            &\log(2+wy)(2+wy)^{\sig-3}
            e^{\ii\Phi_+(r,y,w)}
            \\
            &-
            \log2\,2^{\sig-3}
            e^{\ii\Phi_+(r,y,0)}
        \bigg]
        \,\dd y.
\end{aligned}
\]

For one boundary mode \(m\), the \(J\)-side compact contribution is grouped as
\[
        G_m^\chi(\s)
        =
        \frac{w}{4\pi^2m^2}
        \left[
            c_+(\s)D_+^\chi(\sig,r_m,w)
            +
            c_-(\s)D_-^\chi(\sig,r_m,w)
        \right],
\]
where
\[
        c_+(\s)=\s(\s-2),
        \qquad
        c_-(\s)=(1-\s)(\s+1).
\]
The derivative is
\[
\begin{aligned}
        \partial_\sig G_m^\chi(\s)
        =
        \frac{w}{4\pi^2m^2}
        \big[
            &c_+'(\s)D_+^\chi
            +
            c_+(\s)\partial_\sig D_+^\chi
            \\
            &+
            c_-'(\s)D_-^\chi
            +
            c_-(\s)\partial_\sig D_-^\chi
        \big],
\end{aligned}
\]
with
\[
        c_+'(\s)=2\s-2,
        \qquad
        c_-'(\s)=-2\s.
\]

The boundary certifier divides the boundary sector into six hooks:
\[
        \text{finite compact boundary},
\]
\[
        \text{midrange compact boundary bridge},
\]
\[
        \text{high-\(\tauh\) compact Taylor boundary},
\]
\[
        \text{compact-complement tail},
\]
\[
        \text{compact-complement nonstationary transition},
\]
and
\[
        \text{compact-complement stationary overlap}.
\]
Each hook produces its own JSON certificate.  These hook certificates are then
summed by the strict boundary merger to form
\[
        \code{bd\_sector.safe.json}.
\]

The finite compact-boundary hook covers the initial bounded range
\[
        10\le\tauh\le25.
\]
It evaluates the compact integrals directly by Arb interval boxes in the
variables
\[
        \sig,\quad \tauh,\quad y,
\]
and includes both the value channel and the \(\sig\)-derivative channel.  The
midrange bridge covers
\[
        25\le\tauh\le65536
\]
by dyadic finite-\(\tauh\) blocks.  Each block uses the same compact-boundary
certifier, and the midrange runner sums the safety-normalized block constants.
The high-\(\tauh\) compact Taylor hook covers
\[
        \tauh\ge65536,
        \qquad
        w\le\frac1{256}.
\]
It removes the apparent singularity at \(w=0\) using the identity
\[
        \frac{f(w)-f(0)}{w}
        =
        \int_0^1
        \partial_w f(\lambda w)
        \,\dd\lambda.
\]
This gives direct interval enclosures for
\[
        H_\pm=\frac{D_\pm^\chi}{w}
\]
on boxes containing \(w=0\).  The \(\sig\)-derivative channel is handled by
differentiating the amplitudes while retaining the same Taylor-in-\(w\)
regularization.

The compact cutoff complement
\[
        1-\chi_Y(y)
\]
is not discarded.  It is certified by three hooks.  First, the large-\(|r|\)
compact-complement tail hook factors out the linear oscillation and bounds
\[
        \int_Y^\infty
        e^{\pm 2\pi\ii r y}H_\pm(\sig,y,w)
        \,\dd y
\]
by two integrations by parts.  It uses derivative envelopes for
\[
        H_\pm',
        \qquad
        H_\pm'',
\]
and for their \(\sig\)-derivatives.  The sigma-aware envelope certificate
records the four quantities
\[
        M_{1},
        \qquad
        M_{2},
        \qquad
        M_{1,\sig},
        \qquad
        M_{2,\sig},
\]
which bound
\[
        |H_\pm'|,
        \quad
        |H_\pm''|,
        \quad
        |\partial_\sig H_\pm'|,
        \quad
        |\partial_\sig H_\pm''|.
\]
These are then converted into a large-\(|r|\) tail constant.

Second, the compact-complement transition region
\[
        Y\le y\le2Y
\]
is split according to the phase derivative threshold
\[
        \lambda_0=0.25.
\]
On boxes where
\[
        |\partial_y\Phi_\pm|\ge\lambda_0,
\]
the nonstationary hook applies integration by parts in \(y\).  The
\(\sig\)-derivative certificate uses the same integration-by-parts operator
with the amplitude replaced by its \(\sig\)-derivative.  Boxes where the
interval enclosure of \(\partial_y\Phi_\pm\) intersects
\[
        [-\lambda_0,\lambda_0]
\]
are not counted in the nonstationary hook.  They are exported to the
stationary-overlap hook.

Third, the stationary-overlap hook certifies exactly the exported boxes by
direct absolute Arb interval bounds.  Thus the transition split is exhaustive:
the phase-separated boxes are counted by the nonstationary hook, and the boxes
not separated from zero are counted by the stationary-overlap hook.  In the
final run the two transition hooks recorded complementary counts:
\[
        485976
\]
phase-separated boxes and
\[
        11688
\]
stationary-overlap boxes.  This verifies that the transition grid was split
rather than partially omitted.

The final boundary sector constants in the strict merged certificate are
\[
        C_{{\rm bd},0}
        =
        185599139689.54844424576273899129572044225947553844933704222323532721371511756646,
\]
and
\[
        C_{{\rm bd},1}
        =
        151534184741.81053762160606677355438429846417653094438319187408516674070499261542.
\]
These constants are the safety-normalized sum of all six boundary hooks.  The
boundary certificate is therefore not a single monolithic numerical estimate.
It is a dependency ledger recording the compact finite range, the dyadic
midrange bridge, the high-\(\tauh\) Taylor range, the large-\(|r|\) complement
tail, the nonstationary transition boxes, and the stationary-overlap transition
boxes.  The final boundary JSON records all hook proof hashes and is included
as the sector
\[
        \JSON{bd}
\]
in the final theorem-level certificate.  The final repository reproduces this
boundary sector through the boundary hook scripts and the strict boundary
merger. 

We summarize the boundary certification as follows.

\begin{theorem}[Certified boundary sector]
On the strip
\[
        \sig\in[0.49,0.51],
        \qquad
        \tauh\ge10,
\]
the boundary residual satisfies
\[
        |\Rbd(\s)|
        \le
        C_{{\rm bd},0}
        \frac{\log\tauh}{\tauh^{3/2}},
\]
and
\[
        |\partial_\sig\Rbd(\s)|
        \le
        C_{{\rm bd},1}
        \frac{\log\tauh}{\tauh^{3/2}},
\]
where
\[
        C_{{\rm bd},0}
        =
        185599139689.54844424576273899129572044225947553844933704222323532721371511756646,
\]
and
\[
        C_{{\rm bd},1}
        =
        151534184741.81053762160606677355438429846417653094438319187408516674070499261542.
\]
The certificate includes the projection
\[
        \Rbd=\PiQ\Gbd
\]
and the derivative projection term
\[
        \partial_\sig\Rbd
        =
        (\partial_\sig\PiQ)\Gbd+\PiQ\partial_\sig\Gbd.
\]
\end{theorem}

\begin{proof}
The boundary region is decomposed into the six hooks described above.  The
finite compact-boundary hook covers \(10\le\tauh\le25\).  The midrange bridge
covers \(25\le\tauh\le65536\).  The Taylor-in-\(w\) hook covers
\(\tauh\ge65536\).  The compact-complement tail hook covers the large-\(|r|\)
tail of \(1-\chi_Y\).  The transition region \(Y\le y\le2Y\) is partitioned
into phase-separated boxes and stationary-overlap boxes.  The former are
certified by integration by parts and the latter by direct absolute enclosure.
Each hook produces both a value-channel and a \(\sig\)-derivative-channel
constant.  The strict boundary merger sums the safety-normalized constants and
records the hook proof hashes.  Since each boundary hook is declared
\(J\)-side projected, the quotient projection and its \(\sig\)-derivative are
included.  This proves the stated boundary-sector bounds.
\end{proof}



% ============================================================
% Step 5: Band, endpoint, and core sectors
% ============================================================

\section{Band, endpoint, and core sector certificates}

We now describe the three sectors whose certification is primarily algebraic or
finite-dimensional after the structural decomposition has been made:
\[
        \Rband,
        \qquad
        \Rend,
        \qquad
        \Rcore.
\]
These sectors play different roles in the quotient expansion.  The band sector
is a quotient-side finite stationary band.  The endpoint sector is a
\(J\)-side projected sector arising from the first cell, endpoint extraction,
and the regularized endpoint scalar \(\Xi_{\rm end}\).  The core sector is a
mixed algebraic sector containing the explicit quotient coefficient mismatch.

The final certificates for these three sectors are:
\[
        \code{band\_sector.safe.json},
        \qquad
        \code{end\_sector.safe.json},
        \qquad
        \code{core\_sector.safe.json}.
\]
Their normalization sides are recorded respectively as
\[
        \JSON{Q},
        \qquad
        \JSON{J\_projected},
        \qquad
        \JSON{mixed}.
\]
Thus the band sector is not multiplied by \(\PiQ\) again, the endpoint sector
is projected by \(\PiQ\), and the core sector records its internal
normalization subledger.

\subsection{The quotient band sector}

Let
\[
        L_0=\frac{\tauh}{2\pi},
\]
and define the two band cutoffs
\[
        N(\tauh)
        =
        \left\lfloor\frac{\tauh}{4\pi}\right\rfloor,
        \qquad
        M(\tauh)
        =
        \left\lfloor\frac{\tauh}{2\pi}\right\rfloor.
\]
The quotient band is the finite interval
\[
        N(\tauh)<m\le M(\tauh).
\]
It is visible to the quotient stationary expansion but is not part of the
corresponding \(J\)-tail stationary packet.  This mismatch produces one of the
explicit terms in the effective bracket.

Define the finite band sums
\[
        S_1^{\rm band}(\s,\tauh)
        =
        \sum_{N(\tauh)<m\le M(\tauh)} m^{-\s},
\]
and
\[
        S_2^{\rm band}(\s,\tauh)
        =
        \sum_{N(\tauh)<m\le M(\tauh)} m^{-(1-\s)}.
\]
The quotient stationary coefficients are
\[
        C_A(\s,\tauh)
        =
        -\ii(2\pi)^{1/2-\sig}
        \tauh^{\sig-3/2}
        e^{\ii\Theta(\tauh)-\ii\pi/4},
\]
and
\[
        C_B(\s,\tauh)
        =
        \ii(2\pi)^{\sig-1/2}
        \tauh^{-\sig-1/2}
        e^{-\ii\Theta(\tauh)+\ii\pi/4},
\]
where
\[
        \Theta(\tauh)
        =
        \tauh\log\left(\frac{\tauh}{2\pi}\right)-\tauh.
\]
The exact band contribution is
\[
        \mathcal E_{\rm band}^{\rm exact}(\s,\tauh)
        =
        C_B(\s,\tauh)S_2^{\rm band}(\s,\tauh)
        +
        \frac{\s-1}{\s}
        C_A(\s,\tauh)S_1^{\rm band}(\s,\tauh).
\]
The extracted main band term is
\[
        \mathcal E_{\rm band}^{\rm main}(\s,\tauh)
        =
        \frac{\ii}{2\tauh}
        \left(
            2^{-\sig}e^{-\ii L}
            -
            2^{\sig-1}e^{\ii L}
        \right),
        \qquad
        L=\tauh\log2.
\]
Therefore the certified band residual is
\[
        \Rband(\s)
        =
        \mathcal E_{\rm band}^{\rm exact}(\s,\tauh)
        -
        \mathcal E_{\rm band}^{\rm main}(\s,\tauh).
\]
This is already in quotient normalization.  The band certificate is therefore
declared with
\[
        \JSON{side = Q}.
\]
No additional factor of \(\PiQ\) is applied.

The \(\sig\)-derivative of the exact band term is explicit.  Since the cutoffs
\(N(\tauh)\) and \(M(\tauh)\) depend on \(\tauh\), not on \(\sig\), we have
\[
        \partial_\sig S_1^{\rm band}
        =
        -\sum_{N<m\le M}(\log m)m^{-\s},
\]
and
\[
        \partial_\sig S_2^{\rm band}
        =
        \sum_{N<m\le M}(\log m)m^{-(1-\s)}.
\]
Also
\[
        \partial_\sig C_A
        =
        \log\left(\frac{\tauh}{2\pi}\right)C_A,
\]
and
\[
        \partial_\sig C_B
        =
        \log\left(\frac{2\pi}{\tauh}\right)C_B.
\]
Thus
\[
\begin{aligned}
        \partial_\sig\mathcal E_{\rm band}^{\rm exact}
        =
        &\,
        (\partial_\sig C_B)S_2^{\rm band}
        +
        C_B\partial_\sig S_2^{\rm band}
        \\
        &+
        \frac{1}{\s^2} C_A S_1^{\rm band}
        +
        \frac{\s-1}{\s}
        (\partial_\sig C_A)S_1^{\rm band}
        +
        \frac{\s-1}{\s}
        C_A\partial_\sig S_1^{\rm band}.
\end{aligned}
\]
The derivative of the extracted main term is
\[
        \partial_\sig\mathcal E_{\rm band}^{\rm main}
        =
        -\frac{\ii\log2}{2\tauh}
        \left(
            2^{-\sig}e^{-\ii L}
            +
            2^{\sig-1}e^{\ii L}
        \right).
\]
Therefore
\[
        \partial_\sig\Rband
        =
        \partial_\sig\mathcal E_{\rm band}^{\rm exact}
        -
        \partial_\sig\mathcal E_{\rm band}^{\rm main}.
\]

In the certificate, the band residual is bounded by a conservative Poisson
endpoint envelope.  The envelope accounts for the two dual stationary endpoint
contributions, the nonstationary dual-frequency tails, the endpoint
discretization from the floor cutoffs, and the derivative terms above.  The
final safety-normalized constants recorded in the band certificate are
\[
        C_{{\rm band},0}
        =
        5100100.000005100000,
\]
and
\[
        C_{{\rm band},1}
        =
        9400100.000009400000.
\]
Thus
\[
        |\Rband(\s)|
        \le
        C_{{\rm band},0}
        \frac{\log\tauh}{\tauh^{3/2}},
\]
and
\[
        |\partial_\sig\Rband(\s)|
        \le
        C_{{\rm band},1}
        \frac{\log\tauh}{\tauh^{3/2}}.
\]

\subsection{The endpoint sector}

The endpoint sector contains the first-cell contribution, the leading endpoint
extraction, the \(\zeta(2)\)-cancellation, the regularized
\(\Xi_{\rm end}\)-layer, and the post-integration-by-parts remainder.  It is a
\(J\)-side sector and is therefore declared as
\[
        \JSON{side = J\_projected}.
\]
The projected endpoint residual is
\[
        \Rend(\s)
        =
        \PiQ(\s)G_{\rm end}^{J}(\s),
\]
where
\[
        G_{\rm end}^{J}
        =
        J_1
        +
        E_{\rm cancel}^{(0)}
        +
        \left(E_\Xi-\mathcal L_{\rm end}^{(4)}\right)
        +
        E_{\rm end}^{\rm rem}.
\]
Consequently,
\[
        \partial_\sig\Rend(\s)
        =
        \left(\partial_\sig\PiQ(\s)\right)G_{\rm end}^{J}(\s)
        +
        \PiQ(\s)\partial_\sig G_{\rm end}^{J}(\s).
\]
The derivative of \(\PiQ\) is included in the endpoint certificate.

The first-cell term is
\[
        J_1(\s)
        =
        \int_1^2
        P(x-1)
        \left[
            \s(\s-2)x^{\s-3}
            +
            (1-\s)(\s+1)x^{-\s-2}
        \right]
        \,\dd x,
\]
where
\[
        P(t)=\frac{t^2}{2}-\frac{t}{2}+\frac{1}{12}.
\]
Its \(\sig\)-derivative is obtained directly:
\[
\begin{aligned}
        \partial_\sig J_1(\s)
        =
        \int_1^2
        P(x-1)
        \big[
            &(2\s-2)x^{\s-3}
            +
            \s(\s-2)(\log x)x^{\s-3}
            \\
            &-
            2\s x^{-\s-2}
            -
            (1-\s)(\s+1)(\log x)x^{-\s-2}
        \big]
        \,\dd x.
\end{aligned}
\]

The regularized endpoint scalar is denoted by
\[
        \Xi_{\rm end}(\tauh).
\]
The extracted endpoint layer is
\[
        \mathcal L_{\rm end}^{(4)}(\s,\tauh)
        =
        -\ii\Xi_{\rm end}(\tauh)
        \left(
            2^{\sig-2}e^{\ii L}
            +
            2^{-\sig-1}e^{-\ii L}
        \right).
\]
The correction
\[
        E_\Xi-\mathcal L_{\rm end}^{(4)}
\]
is purely algebraic.  Its derivative is computed by differentiating the
coefficients and the powers of \(2^\sig\); the scalar
\(\Xi_{\rm end}(\tauh)\) depends on \(\tauh\), not on \(\sig\).

The remaining endpoint term is the post-integration-by-parts remainder
\[
        E_{\rm end}^{\rm rem}.
\]
For the plus family, the remainder integrals have the form
\[
        R_{m,+}^{\eps,{\rm end}}(\s)
        =
        \frac{1}{\ii}
        \int_2^\infty
        e^{\ii\Phi_{m,+}^{\eps}(x)}
        \left(
            \frac{x^{\sig-3}}
            {(\Phi_{m,+}^{\eps})'(x)}
        \right)'
        \,\dd x.
\]
For the minus family,
\[
        R_{m,-}^{\eps,{\rm end}}(\s)
        =
        \frac{1}{\ii}
        \int_2^\infty
        e^{\ii\Phi_{m,-}^{\eps}(x)}
        \left(
            \frac{x^{-\sig-2}}
            {(\Phi_{m,-}^{\eps})'(x)}
        \right)'
        \,\dd x.
\]
The \(\sig\)-derivatives are obtained by inserting the corresponding
\(\log x\) factors into the amplitudes.  The endpoint certificate bounds these
remainders by conservative integration-by-parts envelopes.

The endpoint sector certificate records three components:
\[
        \text{first-cell plus }\zeta(2)\text{ cancellation},
\]
\[
        \Xi_{\rm end}\text{ coefficient correction},
\]
and
\[
        \text{post-IBP endpoint remainder}.
\]
The final safety-normalized endpoint constants are
\[
        C_{{\rm end},0}
        =
        71000100.000071000000,
\]
and
\[
        C_{{\rm end},1}
        =
        281000100.000281000000.
\]
Therefore
\[
        |\Rend(\s)|
        \le
        C_{{\rm end},0}
        \frac{\log\tauh}{\tauh^{3/2}},
\]
and
\[
        |\partial_\sig\Rend(\s)|
        \le
        C_{{\rm end},1}
        \frac{\log\tauh}{\tauh^{3/2}}.
\]

\subsection{The core sector}

The core sector contains algebraic corrections left after the main quotient
matching has been performed.  It is marked as
\[
        \JSON{side = mixed}
\]
because its subcomponents may have different normalizations.  The explicit
component certified in the final run is the quotient-side coefficient mismatch.

The quotient stationary coefficient ratio is
\[
        \rho_Q(\s,\tauh)
        =
        \frac{\s-1}{\s}
        \frac{C_A(\s,\tauh)}{C_B(\s,\tauh)}.
\]
The corresponding \(J\)-packet ratio is
\[
        \rho_J(\s,\tauh)
        =
        -\ii
        \frac{\s(\s-2)}{(1-\s)(\s+1)}
        \left(\frac{2\pi}{\tauh}\right)^{1-2\sig}
        e^{2\ii\Theta(\tauh)}.
\]
The exact mismatch is
\[
        \Delta_\rho(\s,\tauh)
        =
        \rho_Q(\s,\tauh)-\rho_J(\s,\tauh).
\]
A direct simplification gives
\[
        \Delta_\rho(\s,\tauh)
        =
        -\ii
        \left(\frac{2\pi}{\tauh}\right)^{1-2\sig}
        e^{2\ii\Theta(\tauh)}
        \frac{\s^2-\s+1}{\s(1-\s)(\s+1)}.
\]
The corresponding core contribution is
\[
        G_{\rho}^{\rm core}(\s)
        =
        C_B(\s,\tauh)
        \Delta_\rho(\s,\tauh)
        \sum_{m\le N(\tauh)}m^{-\s},
\]
with
\[
        N(\tauh)
        =
        \left\lfloor \frac{\tauh}{4\pi}\right\rfloor.
\]
This term is already in quotient normalization.  It must not be multiplied by
\(\PiQ\) again.

The derivative is
\[
\begin{aligned}
        \partial_\sig G_{\rho}^{\rm core}
        =
        &(\partial_\sig C_B)\Delta_\rho
        \sum_{m\le N}m^{-\s}
        \\
        &+
        C_B(\partial_\sig\Delta_\rho)
        \sum_{m\le N}m^{-\s}
        \\
        &+
        C_B\Delta_\rho
        \partial_\sig
        \sum_{m\le N}m^{-\s}.
\end{aligned}
\]
Here
\[
        \partial_\sig C_B
        =
        \log\left(\frac{2\pi}{\tauh}\right)C_B,
\]
and
\[
        \partial_\sig
        \sum_{m\le N}m^{-\s}
        =
        -\sum_{m\le N}(\log m)m^{-\s}.
\]
Writing
\[
        Q(\s)
        =
        \frac{\s^2-\s+1}{\s(1-\s)(\s+1)},
\]
one may compute
\[
        \partial_\sig\Delta_\rho
        =
        -\ii
        \left(\frac{2\pi}{\tauh}\right)^{1-2\sig}
        e^{2\ii\Theta}
        \left[
            -2\log\left(\frac{2\pi}{\tauh}\right)Q(\s)
            +
            Q'(\s)
        \right].
\]
The certificate bounds the Dirichlet polynomial and its derivative by
elementary integral estimates.  Since this component is \(Q\)-side, the core
certificate records
\[
        \JSON{projection = none}
\]
for this subcomponent.

The final safety-normalized core constants are
\[
        C_{{\rm core},0}
        =
        100.15143013245947061906516300167251169542246109728103520969123297128614893010207,
\]
and
\[
        C_{{\rm core},1}
        =
        101.72541371865680722623693802471509402877474481545266923344493309421933909293843.
\]
Thus
\[
        |\Rcore(\s)|
        \le
        C_{{\rm core},0}
        \frac{\log\tauh}{\tauh^{3/2}},
\]
and
\[
        |\partial_\sig\Rcore(\s)|
        \le
        C_{{\rm core},1}
        \frac{\log\tauh}{\tauh^{3/2}}.
\]

Combining these three subsectors gives the band, endpoint, and core entries of
the final certificate.  Their constants are included in the final sums
\[
        \Czero=\sum_\nu C_{\nu,0}^{\rm safe},
        \qquad
        \Cone=\sum_\nu C_{\nu,1}^{\rm safe}.
\]
The normalization ledger ensures that the band and core quotient-side pieces
are not projected twice, while the endpoint sector is projected exactly once
and includes the derivative of the projection.

\begin{theorem}[Certified band, endpoint, and core sectors]
On
\[
        \sig\in[0.49,0.51],
        \qquad
        \tauh\ge10,
\]
the three sectors satisfy
\[
        |\Rband(\s)|
        \le
        C_{{\rm band},0}
        \frac{\log\tauh}{\tauh^{3/2}},
        \qquad
        |\partial_\sig\Rband(\s)|
        \le
        C_{{\rm band},1}
        \frac{\log\tauh}{\tauh^{3/2}},
\]
with
\[
        C_{{\rm band},0}=5100100.000005100000,
        \qquad
        C_{{\rm band},1}=9400100.000009400000;
\]
\[
        |\Rend(\s)|
        \le
        C_{{\rm end},0}
        \frac{\log\tauh}{\tauh^{3/2}},
        \qquad
        |\partial_\sig\Rend(\s)|
        \le
        C_{{\rm end},1}
        \frac{\log\tauh}{\tauh^{3/2}},
\]
with
\[
        C_{{\rm end},0}=71000100.000071000000,
        \qquad
        C_{{\rm end},1}=281000100.000281000000;
\]
and
\[
        |\Rcore(\s)|
        \le
        C_{{\rm core},0}
        \frac{\log\tauh}{\tauh^{3/2}},
        \qquad
        |\partial_\sig\Rcore(\s)|
        \le
        C_{{\rm core},1}
        \frac{\log\tauh}{\tauh^{3/2}},
\]
with
\[
        C_{{\rm core},0}
        =
        100.15143013245947061906516300167251169542246109728103520969123297128614893010207,
\]
and
\[
        C_{{\rm core},1}
        =
        101.72541371865680722623693802471509402877474481545266923344493309421933909293843.
\]
\end{theorem}

\begin{proof}
The band sector is certified directly in quotient normalization by bounding
the exact finite band residual after subtracting its endpoint main term.  The
endpoint sector is first bounded on the \(J\)-side and then projected by
\(\PiQ\), with the derivative term
\((\partial_\sig\PiQ)G_{\rm end}^{J}\) included.  The core sector records the
explicit quotient-side coefficient mismatch and therefore applies no further
projection to that subcomponent.  Each sector certificate contains the value
constant, the \(\sig\)-derivative constant, a coverage hash, a normalization
hash, and a proof hash.  The displayed constants are the safety-normalized
constants accepted by the final merger.
\end{proof}


% ============================================================
% Step 6: Far stationary and global nonstationary sectors
% ============================================================

\section{Far stationary and global nonstationary sector certificates}

We now describe the two remaining analytic tail sectors:
\[
        \Rfar
        \qquad\text{and}\qquad
        \Rnonstat.
\]
Both sectors are naturally produced on the \(J\)-side and are therefore
declared in the certificate with
\[
        \JSON{side = J\_projected}.
\]
Consequently, both sectors are projected into quotient normalization by
\[
        \PiQ(\s)
        =
        \frac{\ii\tauh}{(1-\s)(\s+1)}.
\]
If
\[
        G_{\rm far}^{J}(\s)
        \quad\text{and}\quad
        G_{\rm nonstat}^{J}(\s)
\]
denote the corresponding \(J\)-side residuals, then
\[
        \Rfar(\s)=\PiQ(\s)G_{\rm far}^{J}(\s),
\]
and
\[
        \Rnonstat(\s)=\PiQ(\s)G_{\rm nonstat}^{J}(\s).
\]
The derivative certificates therefore include the projection derivative:
\[
        \partial_\sig\Rfar(\s)
        =
        (\partial_\sig\PiQ(\s))G_{\rm far}^{J}(\s)
        +
        \PiQ(\s)\partial_\sig G_{\rm far}^{J}(\s),
\]
and
\[
        \partial_\sig\Rnonstat(\s)
        =
        (\partial_\sig\PiQ(\s))G_{\rm nonstat}^{J}(\s)
        +
        \PiQ(\s)\partial_\sig G_{\rm nonstat}^{J}(\s),
\]
where
\[
        \partial_\sig\PiQ(\s)
        =
        \PiQ(\s)
        \left(
            \frac{1}{1-\s}
            -
            \frac{1}{\s+1}
        \right).
\]
This projection derivative is part of the normalization contract and is
recorded by the sector normalization hashes.

\subsection{The far stationary sector}

The far stationary sector accounts for the stationary-family residuals that
remain after subtracting the leading stationary packets away from the boundary
transition near \(x=2\).  The relevant stationary families are
\[
        I_{m,+}^{-}(\s)
        =
        \int_2^\infty
        x^{\sig-3}
        e^{\ii(-2\pi m x+\tauh\log x)}
        \,\dd x,
\]
and
\[
        I_{m,-}^{+}(\s)
        =
        \int_2^\infty
        x^{-\sig-2}
        e^{\ii(2\pi m x-\tauh\log x)}
        \,\dd x.
\]
Both have stationary point
\[
        x_{m,*}
        =
        \frac{\tauh}{2\pi m}.
\]
This stationary point lies in the \(J\)-tail \(x\ge2\) precisely when
\[
        m\le \frac{\tauh}{4\pi}.
\]

As in the boundary analysis, set
\[
        \alpha=\frac{\tauh}{4\pi},
        \qquad
        w=\tauh^{-1/2},
        \qquad
        r_m=w(m-\alpha).
\]
The boundary transition corresponds to \(r_m\) bounded.  The far stationary
sector is separated from the boundary by imposing
\[
        r_m\le -R_{\rm far},
\]
with the certified reference value
\[
        R_{\rm far}=4.
\]
Equivalently, the far stationary mode set is
\[
        \mathcal M_{\rm far}(\tauh)
        =
        \left\{
            m\in\mathbb N:
            m\le
            \frac{\tauh}{4\pi}
            -
            R_{\rm far}\sqrt{\tauh}
        \right\}.
\]
This separation ensures that the stationary point is uniformly away from the
endpoint \(x=2\).  The endpoint-stationary transition window is not included in
this sector; it belongs to the boundary sector.

The leading stationary phase is
\[
        \Phi_m(\tauh)
        =
        \tauh\log\left(\frac{\tauh}{2\pi m}\right)-\tauh.
\]
For the plus-amplitude stationary family, the leading stationary packet is
\[
        \operatorname{SP}_{m,+}(\s)
        =
        (2\pi)^{5/2-\sig}
        \tauh^{\sig-5/2}
        m^{2-\sig}
        e^{\ii\Phi_m(\tauh)-\ii\pi/4}.
\]
For the minus-amplitude stationary family, the leading stationary packet is
\[
        \operatorname{SP}_{m,-}(\s)
        =
        (2\pi)^{\sig+3/2}
        \tauh^{-\sig-3/2}
        m^{\sig+1}
        e^{-\ii\Phi_m(\tauh)+\ii\pi/4}.
\]
The exact far stationary remainders are defined by
\[
        D_{m,+}^{\rm far}(\s)
        =
        I_{m,+}^{-}(\s)
        -
        \operatorname{SP}_{m,+}(\s),
\]
and
\[
        D_{m,-}^{\rm far}(\s)
        =
        I_{m,-}^{+}(\s)
        -
        \operatorname{SP}_{m,-}(\s).
\]
This is the object certified by the far sector: the exact integral minus the
explicit leading stationary packet.  No asymptotic symbol is part of the
definition of the residual.

The \(J\)-side far residual is
\[
        G_{\rm far}^{J}(\s)
        =
        \sum_{m\in\mathcal M_{\rm far}(\tauh)}
        \frac{1}{4\pi^2m^2}
        \left[
            c_+(\s)D_{m,+}^{\rm far}(\s)
            +
            c_-(\s)D_{m,-}^{\rm far}(\s)
        \right],
\]
where
\[
        c_+(\s)=\s(\s-2),
        \qquad
        c_-(\s)=(1-\s)(\s+1).
\]
The derivative is
\[
\begin{aligned}
        \partial_\sig G_{\rm far}^{J}(\s)
        =
        \sum_{m\in\mathcal M_{\rm far}(\tauh)}
        \frac{1}{4\pi^2m^2}
        \big[
            &c_+'(\s)D_{m,+}^{\rm far}
            +
            c_+(\s)\partial_\sig D_{m,+}^{\rm far}
            \\
            &+
            c_-'(\s)D_{m,-}^{\rm far}
            +
            c_-(\s)\partial_\sig D_{m,-}^{\rm far}
        \big],
\end{aligned}
\]
with
\[
        c_+'(\s)=2\s-2,
        \qquad
        c_-'(\s)=-2\s.
\]

The \(\sig\)-derivatives of the exact stationary remainders are also exact.
Since the phases are independent of \(\sig\),
\[
        \partial_\sig I_{m,+}^{-}(\s)
        =
        \int_2^\infty
        (\log x)x^{\sig-3}
        e^{\ii(-2\pi m x+\tauh\log x)}
        \,\dd x,
\]
and
\[
        \partial_\sig I_{m,-}^{+}(\s)
        =
        -
        \int_2^\infty
        (\log x)x^{-\sig-2}
        e^{\ii(2\pi m x-\tauh\log x)}
        \,\dd x.
\]
The stationary packets satisfy
\[
        \partial_\sig\operatorname{SP}_{m,+}(\s)
        =
        \log\left(\frac{\tauh}{2\pi m}\right)
        \operatorname{SP}_{m,+}(\s),
\]
and
\[
        \partial_\sig\operatorname{SP}_{m,-}(\s)
        =
        -
        \log\left(\frac{\tauh}{2\pi m}\right)
        \operatorname{SP}_{m,-}(\s).
\]
Thus
\[
        \partial_\sig D_{m,+}^{\rm far}
        =
        \partial_\sig I_{m,+}^{-}
        -
        \log\left(\frac{\tauh}{2\pi m}\right)
        \operatorname{SP}_{m,+},
\]
and
\[
        \partial_\sig D_{m,-}^{\rm far}
        =
        \partial_\sig I_{m,-}^{+}
        +
        \log\left(\frac{\tauh}{2\pi m}\right)
        \operatorname{SP}_{m,-}.
\]

For uniform certification, the far verifier uses scaled stationary variables.
Let
\[
        u=\frac{x}{x_{m,*}},
        \qquad
        x_{m,*}=\frac{\tauh}{2\pi m}.
\]
Define
\[
        \psi(u)=\log u-u+1.
\]
Then
\[
        \psi(1)=0,
        \qquad
        \psi'(1)=0,
        \qquad
        \psi''(1)=-1.
\]
The scaled stationary integrals are
\[
        \mathfrak I_+(\sig,\tauh,\beta)
        =
        \int_\beta^\infty
        u^{\sig-3}
        e^{\ii\tauh\psi(u)}
        \,\dd u,
\]
and
\[
        \mathfrak I_-(\sig,\tauh,\beta)
        =
        \int_\beta^\infty
        u^{-\sig-2}
        e^{-\ii\tauh\psi(u)}
        \,\dd u,
\]
where
\[
        \beta=\frac{4\pi m}{\tauh}.
\]
The leading scaled stationary models are
\[
        \mathfrak S_+(\tauh)
        =
        e^{-\ii\pi/4}\sqrt{\frac{2\pi}{\tauh}},
\]
and
\[
        \mathfrak S_-(\tauh)
        =
        e^{\ii\pi/4}\sqrt{\frac{2\pi}{\tauh}}.
\]
The scaled remainders are
\[
        \mathfrak D_+
        =
        \mathfrak I_+-\mathfrak S_+,
        \qquad
        \mathfrak D_-
        =
        \mathfrak I_- - \mathfrak S_-.
\]
The far-sector envelope certifies these scaled remainders uniformly in the far
stationary range, including their \(\sig\)-derivatives.  The final certificate
uses a deliberately conservative envelope.  The far sector safe constants are
\[
        C_{{\rm far},0}
        =
        50000000100.050000000000,
\]
and
\[
        C_{{\rm far},1}
        =
        100000000100.100000000000.
\]
Therefore
\[
        |\Rfar(\s)|
        \le
        C_{{\rm far},0}
        \frac{\log\tauh}{\tauh^{3/2}},
\]
and
\[
        |\partial_\sig\Rfar(\s)|
        \le
        C_{{\rm far},1}
        \frac{\log\tauh}{\tauh^{3/2}}.
\]

\subsection{The global nonstationary sector}

The global nonstationary sector contains the four-family tail pieces that are
not assigned to the band, boundary, endpoint, far stationary, floor, or core
sectors.  The four tail families are
\[
        I_{m,+}^{\eps}(\s)
        =
        \int_2^\infty
        x^{\sig-3}
        e^{\ii(\eps 2\pi m x+\tauh\log x)}
        \,\dd x,
\]
and
\[
        I_{m,-}^{\eps}(\s)
        =
        \int_2^\infty
        x^{-\sig-2}
        e^{\ii(\eps 2\pi m x-\tauh\log x)}
        \,\dd x,
\]
with
\[
        \eps\in\{+1,-1\}.
\]
The phase derivatives are
\[
        (\Phi_{m,+}^{\eps})'(x)
        =
        \eps 2\pi m+\frac{\tauh}{x},
\]
and
\[
        (\Phi_{m,-}^{\eps})'(x)
        =
        \eps 2\pi m-\frac{\tauh}{x}.
\]
For \(\tauh>0\), the families
\[
        I_{m,+}^{+},
        \qquad
        I_{m,-}^{-}
\]
are always nonstationary.  The families
\[
        I_{m,+}^{-},
        \qquad
        I_{m,-}^{+}
\]
may have stationary points, but only when
\[
        m\le\frac{\tauh}{4\pi}.
\]
The stationary pieces in the far range are assigned to \(\Rfar\), and the
endpoint-stationary transition window is assigned to \(\Rbd\).  The remaining
parts are assigned to the global nonstationary sector.

Let
\[
        \mathscr N
        \subset
        \mathbb N\times\{+,-\}\times\{+,-\}
\]
be the certified nonstationary index set after removing the already assigned
band, boundary, endpoint, far, floor, and core pieces.  For
\[
        \kappa\in\{+,-\},
        \qquad
        \eps\in\{+,-\},
\]
write
\[
        I_{m,\kappa}^{\eps}(\s)
\]
for the corresponding oscillatory integral.  The \(J\)-side nonstationary
residual is
\[
        G_{\rm nonstat}^{J}(\s)
        =
        \frac{1}{4\pi^2}
        \sum_{(m,\kappa,\eps)\in\mathscr N}
        \frac{1}{m^2}
        c_\kappa(\s)
        I_{m,\kappa}^{\eps}(\s),
\]
where
\[
        c_+(\s)=\s(\s-2),
        \qquad
        c_-(\s)=(1-\s)(\s+1).
\]
Its derivative is
\[
\begin{aligned}
        \partial_\sig G_{\rm nonstat}^{J}(\s)
        =
        \frac{1}{4\pi^2}
        \sum_{(m,\kappa,\eps)\in\mathscr N}
        \frac{1}{m^2}
        \left[
            c_\kappa'(\s)
            I_{m,\kappa}^{\eps}(\s)
            +
            c_\kappa(\s)
            \partial_\sig I_{m,\kappa}^{\eps}(\s)
        \right].
\end{aligned}
\]
Since the phases do not depend on \(\sig\),
\[
        \partial_\sig I_{m,+}^{\eps}(\s)
        =
        \int_2^\infty
        (\log x)x^{\sig-3}
        e^{\ii\Phi_{m,+}^{\eps}(x)}
        \,\dd x,
\]
and
\[
        \partial_\sig I_{m,-}^{\eps}(\s)
        =
        -
        \int_2^\infty
        (\log x)x^{-\sig-2}
        e^{\ii\Phi_{m,-}^{\eps}(x)}
        \,\dd x.
\]

The certification is by integration by parts on regions where the phase
derivative is separated from zero.  For a general phase derivative
\[
        p(x)=\Phi'(x),
\]
define
\[
        \mathcal L_p A
        =
        \left(\frac{A}{p}\right)'
        =
        \frac{A'}{p}
        -
        \frac{A p'}{p^2}.
\]
Then
\[
        \int_a^\infty A(x)e^{\ii\Phi(x)}\,\dd x
        =
        -\frac{A(a)e^{\ii\Phi(a)}}{\ii p(a)}
        -
        \frac{1}{\ii}
        \int_a^\infty
        \mathcal L_pA(x)e^{\ii\Phi(x)}
        \,\dd x.
\]
A second integration by parts gives the certified bound
\[
\begin{aligned}
        \left|
        \int_a^\infty A e^{\ii\Phi}\,\dd x
        \right|
        \le
        &
        \frac{|A(a)|}{|p(a)|}
        +
        \frac{|(\mathcal L_pA)(a)|}{|p(a)|}
        \\
        &+
        \int_a^\infty
        |\mathcal L_p^2A(x)|
        \,\dd x.
\end{aligned}
\]
The derivative certificate uses the same operator with
\[
        A
        \quad\text{replaced by}\quad
        \partial_\sig A.
\]
This is valid because the phases are independent of \(\sig\).

The global nonstationary certificate records a conservative envelope for the
always-nonstationary families, the potentially stationary families outside
their stationary ranges, and a partition guard ensuring that no region already
assigned to another sector is counted again.  The final safety-normalized constants are
\[
        C_{{\rm nonstat},0}
        =
        20000000100.020000000000,
\]
and
\[
        C_{{\rm nonstat},1}
        =
        40000000100.040000000000.
\]
Thus
\[
        |\Rnonstat(\s)|
        \le
        C_{{\rm nonstat},0}
        \frac{\log\tauh}{\tauh^{3/2}},
\]
and
\[
        |\partial_\sig\Rnonstat(\s)|
        \le
        C_{{\rm nonstat},1}
        \frac{\log\tauh}{\tauh^{3/2}}.
\]

The final repository
\[
        \texttt{https://github.com/vicblab/RH-inner-strip-certifier}
\]
contains the scripts and tutorial-style instructions used to reproduce these
sector certificates, including the one-command driver
\[
        \code{./run\_all\_certification.sh}.
\]
The repository README records the intended strip
\[
        \sig\in[0.49,0.51],
\]
the bracket parameters
\[
        u_0=0.01,
        \qquad
        g_0=0.2,
\]
the target height
\[
        T=10^{30},
\]
and the expected merged constants
\[
        C_{\rm even}\approx2.5567524018976995\times10^{11},
\]
\[
        C_{\sig}\approx2.9182458524367624\times10^{11},
\]
\[
        C_{\rm inner}\approx3.879837844533457\times10^{11}.
\]
It also states that the full run creates
\[
        \code{INNER\_STRIP\_CERTIFICATE\_SAFE.json}
\]
and that the final JSON must contain
\[
        \JSON{status = inner\_full\_certificate}
\]
and
\[
        \JSON{height\_condition\_holds = true}.
\]
This repository is the public reproduction artifact for the certified
constants quoted in this paper. \cite{InnerRepo}

We summarize the far and global nonstationary sectors as follows.

\begin{theorem}[Certified far and global nonstationary sectors]
On
\[
        \sig\in[0.49,0.51],
        \qquad
        \tauh\ge10,
\]
the far stationary residual satisfies
\[
        |\Rfar(\s)|
        \le
        C_{{\rm far},0}
        \frac{\log\tauh}{\tauh^{3/2}},
\]
and
\[
        |\partial_\sig\Rfar(\s)|
        \le
        C_{{\rm far},1}
        \frac{\log\tauh}{\tauh^{3/2}},
\]
with
\[
        C_{{\rm far},0}
        =
        50000000100.050000000000,
\]
and
\[
        C_{{\rm far},1}
        =
        100000000100.100000000000.
\]
The global nonstationary residual satisfies
\[
        |\Rnonstat(\s)|
        \le
        C_{{\rm nonstat},0}
        \frac{\log\tauh}{\tauh^{3/2}},
\]
and
\[
        |\partial_\sig\Rnonstat(\s)|
        \le
        C_{{\rm nonstat},1}
        \frac{\log\tauh}{\tauh^{3/2}},
\]
with
\[
        C_{{\rm nonstat},0}
        =
        20000000100.020000000000,
\]
and
\[
        C_{{\rm nonstat},1}
        =
        40000000100.040000000000.
\]
Both sectors are \(J\)-side projected and include the derivative of the
quotient projection.
\end{theorem}

\begin{proof}
The far sector is the exact stationary-family residual after subtracting the
leading stationary packets in the region \(r_m\le -R_{\rm far}\).  The
stationary point is uniformly separated from the endpoint \(x=2\), and the
scaled stationary variable reduces the certification to uniform bounds for the
scaled remainders and their \(\sig\)-derivatives.  The global nonstationary
sector consists of the remaining four-family tail pieces after removing the
band, boundary, endpoint, far, floor, and core sectors.  On this set the phase
derivative is bounded away from zero, so repeated integration by parts gives
uniform value and \(\sig\)-derivative bounds.  Both \(J\)-side residuals are
then projected by \(\PiQ\), and the derivative certificates include
\((\partial_\sig\PiQ)G^{J}\).  The displayed constants are the
safety-normalized constants recorded in the final sector JSON files and
accepted by the final strict merger.
\end{proof}


% ============================================================
% Step 7: Final merge, height computation, and theorem assembly
% ============================================================

\section{Final merge and height computation}

We now assemble the seven residual-sector certificates and the finite-width
bracket certificate into the final inner-strip certificate.  The final merger is
the strict bracket-aware verifier
\[
        \code{merge\_inner\_certificate\_with\_bracket\_safe.py}.
\]
It requires the seven residual-sector files
\[
        \code{band\_sector.safe.json},
        \quad
        \code{bd\_sector.safe.json},
        \quad
        \code{far\_sector.safe.json},
        \quad
        \code{end\_sector.safe.json},
\]
\[
        \code{nonstat\_sector.safe.json},
        \quad
        \code{floor\_sector.safe.json},
        \quad
        \code{core\_sector.safe.json},
\]
and it also requires the finite-width bracket certificate
\[
        \code{inner\_bracket\_vector.safe.json}.
\]
The final certificate is not allowed to be produced from the residual-sector
certificates alone.  This is important: the obstruction requires both the
residual upper bound and the finite-width bracket lower bound.

The strict final merger checks the following conditions.

First, every residual-sector file must have
\[
        \JSON{status = proved}.
\]
Second, every residual sector must declare one of the permitted normalization
sides:
\[
        \Qside,
        \qquad
        \Jproj,
        \qquad
        \Zeroside,
        \qquad
        \Mixedside.
\]
Third, every residual sector must contain both safety-normalized constants
\[
        \JSON{C\_value\_upper\_safe}
\]
and
\[
        \JSON{C\_sigma\_derivative\_upper\_safe}.
\]
Fourth, the finite-width bracket certificate must have
\[
        \JSON{status = proved}
\]
and must certify
\[
        \JSON{g0\_bracket\_vector\_lower}\ge0.2.
\]
Only after these conditions are satisfied does the final merger compute the
residual vector constant and check the height inequality.

Let
\[
        C_{\nu,0}^{\rm safe}
\]
and
\[
        C_{\nu,1}^{\rm safe}
\]
denote the value-channel and \(\sig\)-derivative-channel constants for a
sector
\[
        \nu\in
        \{
        {\rm band},
        {\rm bd},
        {\rm far},
        {\rm end},
        {\rm nonstat},
        {\rm floor},
        {\rm core}
        \}.
\]
The final value-channel constant is
\[
        \Czero
        =
        \sum_\nu C_{\nu,0}^{\rm safe}.
\]
The final derivative-channel constant is
\[
        \Cone
        =
        \sum_\nu C_{\nu,1}^{\rm safe}.
\]
The two channels are then combined as a Euclidean vector bound:
\[
        \Cinner
        =
        \sqrt{\Czero^2+\Cone^2}.
\]
In the final certificate, the seven sector safe constants are as follows.

For the band sector,
\[
        C_{{\rm band},0}^{\rm safe}
        =
        5100100.000005100000,
\]
and
\[
        C_{{\rm band},1}^{\rm safe}
        =
        9400100.000009400000.
\]
For the boundary sector,
\[
        C_{{\rm bd},0}^{\rm safe}
        =
        185599139689.54844424576273899129572044225947553844933704222323532721371511756646,
\]
and
\[
        C_{{\rm bd},1}^{\rm safe}
        =
        151534184741.81053762160606677355438429846417653094438319187408516674070499261542.
\]
For the far stationary sector,
\[
        C_{{\rm far},0}^{\rm safe}
        =
        50000000100.050000000000,
\]
and
\[
        C_{{\rm far},1}^{\rm safe}
        =
        100000000100.100000000000.
\]
For the endpoint sector,
\[
        C_{{\rm end},0}^{\rm safe}
        =
        71000100.000071000000,
\]
and
\[
        C_{{\rm end},1}^{\rm safe}
        =
        281000100.000281000000.
\]
For the global nonstationary sector,
\[
        C_{{\rm nonstat},0}^{\rm safe}
        =
        20000000100.020000000000,
\]
and
\[
        C_{{\rm nonstat},1}^{\rm safe}
        =
        40000000100.040000000000.
\]
For the floor sector,
\[
        C_{{\rm floor},0}^{\rm safe}=0,
        \qquad
        C_{{\rm floor},1}^{\rm safe}=0.
\]
For the core sector,
\[
        C_{{\rm core},0}^{\rm safe}
        =
        100.15143013245947061906516300167251169542246109728103520969123297128614893010207,
\]
and
\[
        C_{{\rm core},1}^{\rm safe}
        =
        101.72541371865680722623693802471509402877474481545266923344493309421933909293843.
\]

Summing the value-channel constants gives
\[
\begin{aligned}
        \Czero
        &=
        C_{{\rm band},0}^{\rm safe}
        +
        C_{{\rm bd},0}^{\rm safe}
        +
        C_{{\rm far},0}^{\rm safe}
        +
        C_{{\rm end},0}^{\rm safe}
        \\
        &\quad
        +
        C_{{\rm nonstat},0}^{\rm safe}
        +
        C_{{\rm floor},0}^{\rm safe}
        +
        C_{{\rm core},0}^{\rm safe}
        \\
        &=
        255675240189.76995047822220961036088344393198723387179813950427053690494808885261.
\end{aligned}
\]
Summing the derivative-channel constants gives
\[
\begin{aligned}
        \Cone
        &=
        C_{{\rm band},1}^{\rm safe}
        +
        C_{{\rm bd},1}^{\rm safe}
        +
        C_{{\rm far},1}^{\rm safe}
        +
        C_{{\rm end},1}^{\rm safe}
        \\
        &\quad
        +
        C_{{\rm nonstat},1}^{\rm safe}
        +
        C_{{\rm floor},1}^{\rm safe}
        +
        C_{{\rm core},1}^{\rm safe}
        \\
        &=
        291824585243.67624174026287399979132232317927055971912800732675440018563808683476.
\end{aligned}
\]
Therefore
\[
        \Cinner
        =
        \sqrt{\Czero^2+\Cone^2}
\]
is
\[
        \Cinner
        =
        387983784453.34570013486357394399212036290592556640627403482826815918742177195616.
\]
This value is recorded in the final JSON as
\[
        \JSON{C\_inner\_vector\_upper}.
\]

The bracket certificate gives
\[
        \gzero=0.2.
\]
The final obstruction inequality is
\[
        \sqrt T>
        \frac{\Cinner}{\gzero}\log T.
\]
Substituting the certified values gives
\[
        \frac{\Cinner}{\gzero}
        =
        1939918922266.7285006743178697199606018145296278320313701741413407959371088597808.
\]
Thus the inequality to be checked is
\[
        \sqrt T
        >
        1939918922266.7285006743178697199606018145296278320313701741413407959371088597808
        \log T.
\]

The final certificate checks this inequality at
\[
        T=10^{30}.
\]
At this height,
\[
        \sqrt T=10^{15},
\]
while
\[
        \frac{\Cinner}{\gzero}\log T
        =
        1939918922266.7285006743178697199606018145296278320313701741413407959371088597808
        \cdot
        30\log 10.
\]
The strict merger records
\[
        \JSON{height\_condition\_holds = true}.
\]
Thus the certified height condition holds at the shared target height
\[
        T=10^{30}.
\]

The height threshold can also be solved directly.  Setting
\[
        \sqrt T=
        \frac{\Cinner}{\gzero}\log T
\]
and solving numerically gives
\[
        T_{\rm min}
        \approx
        1.5867647318922528\times10^{28}.
\]
Equivalently,
\[
        \log_{10}T_{\rm min}
        \approx
        28.2005125390955.
\]
Therefore the clean power-of-ten height statement implied by the certificate is
\[
        T\ge10^{29}.
\]
The repository provides the script
\[
        \code{tools/height\_threshold.py}
\]
to reproduce this computation.  The public repository
\[
        \texttt{https://github.com/vicblab/RH-inner-strip-certifier}
\]
also records the intended target strip, the expected merged constants, the
one-command reproduction script, and the final JSON condition
\[
        \JSON{status = inner\_full\_certificate},
        \qquad
        \JSON{height\_condition\_holds = true}.
\]
The README states that running
\[
        \code{./run\_all\_certification.sh}
\]
creates all sector certificates in
\[
        \code{certs\_inner/}
\]
and then creates
\[
        \code{INNER\_STRIP\_CERTIFICATE\_SAFE.json}.
\]
This is the public reproducibility artifact for the final constants used in
this paper. \cite{InnerRepo}
We now state the final certified merge theorem.

\begin{theorem}[Strict final inner-strip certificate]
The strict bracket-aware final merger produces the certificate
\[
        \code{INNER\_STRIP\_CERTIFICATE\_SAFE.json}
\]
with
\[
        \JSON{status = inner\_full\_certificate}
\]
and
\[
        \JSON{height\_condition\_holds = true}.
\]
The certificate records
\[
        \Czero
        =
        255675240189.76995047822220961036088344393198723387179813950427053690494808885261,
\]
\[
        \Cone
        =
        291824585243.67624174026287399979132232317927055971912800732675440018563808683476,
\]
and
\[
        \Cinner
        =
        387983784453.34570013486357394399212036290592556640627403482826815918742177195616.
\]
Together with
\[
        \gzero=0.2,
\]
the height inequality
\[
        \sqrt T>
        \frac{\Cinner}{\gzero}\log T
\]
holds at
\[
        T=10^{30}.
\]
Moreover, the numerical threshold is
\[
        T_{\rm min}
        \approx
        1.5867647318922528\times10^{28},
\]
so the clean certified height statement is
\[
        T\ge10^{29}.
\]
\end{theorem}

\begin{proof}
The final merger first verifies that all seven residual sectors are present and
have status \(\JSON{proved}\).  It then verifies that the finite-width bracket
certificate is present and has status \(\JSON{proved}\), with
\[
        \JSON{g0\_bracket\_vector\_lower}\ge0.2.
\]
It sums the seven value-channel constants to obtain \(\Czero\), sums the seven
\(\sig\)-derivative-channel constants to obtain \(\Cone\), and computes
\[
        \Cinner=\sqrt{\Czero^2+\Cone^2}.
\]
Finally, it evaluates
\[
        \sqrt T>
        \frac{\Cinner}{\gzero}\log T
\]
at
\[
        T=10^{30}.
\]
The computed Boolean field
\[
        \JSON{height\_condition\_holds}
\]
is true.  Therefore the strict final certificate is issued.
\end{proof}

Combining the strict final certificate with the structural quotient obstruction
gives the desired exclusion statement.

\begin{corollary}[Certified central inner-strip exclusion above \(10^{29}\)]
Assume the structural quotient identity and two-channel obstruction theorem
proved in the analytic companion paper.  Then the Riemann zeta function has no
off-critical nontrivial zeros in the central inner strip
\[
        0<|\Re s-\tfrac12|\le0.01
\]
for
\[
        |\Im s|\ge10^{29}.
\]
In particular, the exclusion holds at the target height
\[
        |\Im s|\ge10^{30}.
\]
\end{corollary}

\begin{proof}
At a nontrivial zero, the structural identity gives
\[
        \NQ(s)=0.
\]
The certified expansion gives
\[
        \NQ(s)
        =
        \frac{\ii}{\tauh}\Beff(\sig,\tauh)
        +
        \Rtot(s).
\]
The finite-width bracket certificate gives
\[
        \|\Bvec(u,\tauh)\|\ge\gzero=0.2
\]
for
\[
        0\le u\le0.01.
\]
The residual certificate gives
\[
        \|\Rtot(s)\|
        \le
        \Cinner\frac{\log\tauh}{\tauh^{3/2}}.
\]
If
\[
        \sqrt{\tauh}>
        \frac{\Cinner}{\gzero}\log\tauh,
\]
then the residual is too small to cancel the bracket term.  Hence
\[
        \NQ(s)\ne0,
\]
contradicting the zero assumption.  The certified height computation shows
that this inequality holds for
\[
        \tauh\ge10^{29}.
\]
Thus there are no off-critical zeros in the stated inner strip above that
height.
\end{proof}


\section{Global assembly with the outer-strip certificate}

The certificate proved in this paper is an inner-strip certificate.  Its
certified region is
\[
        \Omega_{\rm in}
        =
        \left\{
            s=\sig+\ii\tauh:
            0<|\sig-\tfrac12|\le0.01,
            \ |\tauh|\ge10^{29}
        \right\}.
\]
It does not by itself certify the full critical strip.  To obtain a global
high-height zero exclusion away from the endpoints \(\sig=0\) and \(\sig=1\),
one must combine it with the outer-strip certificate.

Let \(\delta_{\rm out}>0\) denote the endpoint margin certified by the
outer-strip argument.  Assume that the outer-strip companion theorem \cite{OuterCompanion} proves
that there are no off-critical zeros in
\[
        \Omega_{\rm out}
        =
        \left\{
            s=\sig+\ii\tauh:
            0.01\le\sig\le0.49,
            \ |\tauh|\ge10^{30}
        \right\}
        \cup
        \left\{
            s=\sig+\ii\tauh:
            0.51\le\sig\le0.99,
            \ |\tauh|\ge10^{30}
        \right\}.
\]
Assume also that the outer-strip theorem and the present inner-strip theorem
use compatible quotient-normalization conventions and compatible height
normalizations.

The present paper proves the central part
\[
        0.49\le\sig\le0.51,
        \qquad
        \sig\ne\frac12,
        \qquad
        |\tauh|\ge10^{29}.
\]
Since
\[
        10^{30}\ge10^{29},
\]
the inner-strip certificate is valid throughout the central strip at the
outer-strip target height \(10^{30}\).  Therefore, after combining the outer
and inner certificates, the covered high-height region is
\[
        0.01\le\sig\le1-0.01,
        \qquad
        \sig\ne\frac12,
        \qquad
        |\tauh|\ge10^{30}.
\]

\begin{theorem}[Conditional global high-height assembly]
Assume the structural quotient-obstruction theorem, the finite-width
inner-strip certificate proved in this paper, and the outer-strip certificate \cite{OuterRepo}
on
\[
        [0.01,0.49]\cup[0.51,0.99]
\]
at height \(10^{30}\).  Then there are no off-critical nontrivial zeros of
\(\zeta(s)\) in
\[
        0.01\le\Re s\le1-0.01,
        \qquad
        \Re s\ne\frac12,
        \qquad
        |\Im s|\ge10^{30}.
\]
\end{theorem}

\begin{proof}
The outer-strip certificate excludes off-critical zeros for
\[
        0.01\le\sig\le0.49
        \quad\text{and}\quad
        0.51\le\sig\le0.99
\]
at height \(10^{30}\).  The present inner-strip certificate excludes
off-critical zeros for
\[
        0.49\le\sig\le0.51,
        \qquad
        \sig\ne\frac12,
\]
already for height \(10^{29}\), and hence also for height \(10^{30}\).
The union of these certified regions is
\[
        0.01\le\sig\le0.99,
        \qquad
        \sig\ne\frac12.
\]
Thus no off-critical zero can occur in the stated region at height
\(|\tauh|\ge10^{30}\).
\end{proof}


% ============================================================
% Step 8: Audit appendix, reproducibility protocol, and bibliography
% ============================================================

\appendix

\section{Reproducibility and audit protocol}

This appendix records the exact reproducibility protocol for the inner-strip
certificate.  The public repository for the certificate is
\[
        \texttt{https://github.com/vicblab/RH-inner-strip-certifier}.
\]
The repository README states that the certifier reproduces the inner-strip
residual constants for
\[
        \sig\in[0.49,0.51],
        \qquad
        u_0=0.01,
        \qquad
        g_0=0.2,
        \qquad
        T=10^{30},
\]
and that the final certificate is the file
\[
        \code{INNER\_STRIP\_CERTIFICATE\_SAFE.json}.
\]
It also states that the expected merged constants are approximately
\[
        C_{\rm even}\approx2.5567524018976995\times10^{11},
\]
\[
        C_{\sig}\approx2.9182458524367624\times10^{11},
\]
and
\[
        C_{\rm inner}\approx3.879837844533457\times10^{11}.
\]
The repository provides a one-command reproduction script
\[
        \code{./run\_all\_certification.sh},
\]
which generates all sector certificates in
\[
        \code{certs\_inner/}
\]
and then creates
\[
        \code{INNER\_STRIP\_CERTIFICATE\_SAFE.json}.
\]
The README further specifies that the final JSON must contain
\[
        \JSON{status = inner\_full\_certificate}
\]
and
\[
        \JSON{height\_condition\_holds = true}.
\]
These are the public reproduction conditions for the certificate. \cite{InnerRepo}

The repository is intended to be tested from a clean Python environment.  A
typical run begins with
\[
        \code{python -m venv .venv},
\]
followed by
\[
        \code{source .venv/bin/activate},
\]
\[
        \code{python -m pip install --upgrade pip},
\]
and
\[
        \code{python -m pip install python-flint}.
\]
After installing the dependency, the full certificate is reproduced by running
\[
        \code{./run\_all\_certification.sh}.
\]
This command creates the seven residual-sector certificates, the finite-width
bracket certificate, and the strict final merged certificate.

The complete list of final sector certificates is
\[
        \code{certs\_inner/band\_sector.safe.json},
\]
\[
        \code{certs\_inner/bd\_sector.safe.json},
\]
\[
        \code{certs\_inner/far\_sector.safe.json},
\]
\[
        \code{certs\_inner/end\_sector.safe.json},
\]
\[
        \code{certs\_inner/nonstat\_sector.safe.json},
\]
\[
        \code{certs\_inner/floor\_sector.safe.json},
\]
and
\[
        \code{certs\_inner/core\_sector.safe.json}.
\]
The final strict merger also requires
\[
        \code{certs\_inner/inner\_bracket\_vector.safe.json}.
\]
Thus the minimum complete certificate directory must contain eight files:
seven residual-sector JSON files and one bracket JSON file.

The final theorem-level merger is
\[
        \code{scripts/merge\_inner\_certificate\_with\_bracket\_safe.py}.
\]
It is run as
\[
\begin{aligned}
        &\code{python scripts/merge\_inner\_certificate\_with\_bracket\_safe.py} \\
        &\quad
        \code{--inputs 'certs\_inner/*\_sector.safe.json'} \\
        &\quad
        \code{--bracket certs\_inner/inner\_bracket\_vector.safe.json} \\
        &\quad
        \code{--height 1e30} \\
        &\quad
        \code{--g0 0.2} \\
        &\quad
        \code{--output INNER\_STRIP\_CERTIFICATE\_SAFE.json}.
\end{aligned}
\]
This strict merger checks both the residual certificates and the finite-width
bracket certificate.  It refuses to produce a theorem-level certificate if the
bracket certificate is missing.  This is essential because the residual
constants alone do not prove the obstruction; the finite-width bracket lower
bound is also required.

The final file
\[
        \code{INNER\_STRIP\_CERTIFICATE\_SAFE.json}
\]
must contain
\[
        \JSON{status = inner\_full\_certificate},
\]
\[
        \JSON{height\_condition\_holds = true},
\]
\[
        \JSON{u0 = 0.01},
\]
and
\[
        \JSON{g0\_bracket\_vector\_lower = 0.2}.
\]
It must also contain a
\[
        \JSON{bracket\_certificate}
\]
field recording the file and proof hash of
\[
        \code{certs\_inner/inner\_bracket\_vector.safe.json}.
\]
Finally, it must contain seven sector entries under the field
\[
        \JSON{sectors}.
\]
The seven sector names must be exactly
\[
        \JSON{band},
        \quad
        \JSON{bd},
        \quad
        \JSON{far},
        \quad
        \JSON{end},
        \quad
        \JSON{nonstat},
        \quad
        \JSON{floor},
        \quad
        \JSON{core}.
\]

\section{Sector audit checklist}

Each residual-sector JSON is checked according to the same basic schema.  The
field
\[
        \JSON{status}
\]
must equal
\[
        \JSON{proved}.
\]
The field
\[
        \JSON{side}
\]
must be one of
\[
        \Qside,
        \qquad
        \Jproj,
        \qquad
        \Zeroside,
        \qquad
        \Mixedside.
\]
The fields
\[
        \JSON{C\_value\_upper\_safe}
\]
and
\[
        \JSON{C\_sigma\_derivative\_upper\_safe}
\]
must be present and nonnegative.  The fields
\[
        \JSON{coverage\_hash},
        \qquad
        \JSON{normalization\_hash},
        \qquad
        \JSON{proof\_hash}
\]
must also be present.

The permitted sector sides are fixed by the normalization ledger:
\[
        \JSON{band} \mapsto \Qside,
\]
\[
        \JSON{bd} \mapsto \Jproj,
\]
\[
        \JSON{far} \mapsto \Jproj,
\]
\[
        \JSON{end} \mapsto \Jproj,
\]
\[
        \JSON{nonstat} \mapsto \Jproj,
\]
\[
        \JSON{floor} \mapsto \Zeroside,
\]
and
\[
        \JSON{core} \mapsto \Mixedside.
\]
The band sector is already in quotient normalization and is not projected
again.  The boundary, far, endpoint, and global nonstationary sectors are
\(J\)-side sectors and are projected by \(\PiQ\).  Their derivative
certificates include the term
\[
        (\partial_\sig\PiQ)G^J.
\]
The floor sector is identically zero.  The core sector is mixed because the
explicit coefficient mismatch is quotient-side, while any later algebraic
subcomponents must declare their own normalization.

The boundary sector has its own internal audit.  The file
\[
        \code{bd\_sector.safe.json}
\]
is assembled from the following hooks:
\[
        \text{finite compact boundary},
\]
\[
        \text{midrange compact boundary bridge},
\]
\[
        \text{high-\(\tauh\) compact Taylor boundary},
\]
\[
        \text{compact-complement tail},
\]
\[
        \text{compact-complement nonstationary transition},
\]
and
\[
        \text{compact-complement stationary overlap}.
\]
The boundary sector is accepted only after these hooks are merged into the
single boundary-sector certificate.  The transition hooks record complementary
box counts: phase-separated boxes are assigned to the nonstationary hook, and
boxes whose derivative interval intersects the stationary threshold are
assigned to the stationary-overlap hook.  This prevents the compact-complement
transition region from being silently dropped.

The bracket certificate is checked separately.  It must record
\[
        \JSON{status = proved},
\]
\[
        \JSON{sector = inner\_bracket\_vector},
\]
\[
        \JSON{u0 = 0.01},
\]
and
\[
        \JSON{g0\_bracket\_vector\_lower = 0.2}.
\]
The strict final merger does not accept the earlier critical-line-only bracket
record.  It requires a finite-width bracket certificate.

\section{Expected final constants}

The final JSON produced in the certified run records
\[
        \Czero
        =
        255675240189.76995047822220961036088344393198723387179813950427053690494808885261,
\]
\[
        \Cone
        =
        291824585243.67624174026287399979132232317927055971912800732675440018563808683476,
\]
and
\[
        \Cinner
        =
        387983784453.34570013486357394399212036290592556640627403482826815918742177195616.
\]
The value
\[
        \Cinner
\]
is the Euclidean combination
\[
        \Cinner
        =
        \sqrt{\Czero^2+\Cone^2}.
\]
The final bracket lower bound is
\[
        \gzero=0.2.
\]
Thus
\[
        \frac{\Cinner}{\gzero}
        =
        1939918922266.7285006743178697199606018145296278320313701741413407959371088597808.
\]

The final certificate checks the height inequality
\[
        \sqrt T>
        \frac{\Cinner}{\gzero}\log T
\]
at
\[
        T=10^{30}.
\]
It records
\[
        \JSON{height\_condition\_holds = true}.
\]
Solving the equality
\[
        \sqrt T=
        \frac{\Cinner}{\gzero}\log T
\]
gives
\[
        T_{\rm min}
        \approx
        1.5867647318922528\times10^{28}.
\]
Therefore the clean power-of-ten certified height is
\[
        T\ge10^{29}.
\]
The repository contains the utility
\[
        \code{tools/height\_threshold.py}
\]
to reproduce this calculation.

\section{Hash and dependency policy}

Every certificate JSON contains a proof hash.  The proof hash is not intended
to replace mathematical verification; rather, it prevents accidental
transcription or substitution of JSON files during the merge.  The final JSON
records the proof hash of each of the seven sector certificates and the proof
hash of the bracket certificate.  It also records coverage hashes and
normalization hashes.  These serve as a dependency ledger.

The coverage hash records which part of the decomposition is being certified.
For example, the boundary coverage depends on the compact finite range,
midrange bridge, high-\(\tauh\) Taylor range, compact-complement tail,
nonstationary transition, and stationary-overlap transition.  The normalization
hash records whether the sector is \(Q\)-side, \(J\)-side projected, mixed, or
zero.  In particular, it records whether the quotient projection and its
\(\sig\)-derivative are included.

A verifier or referee should treat the final JSON as valid only if all of the
following hold:
\[
        \JSON{status = inner\_full\_certificate},
\]
\[
        \JSON{height\_condition\_holds = true},
\]
all seven residual sectors are present,
\[
        \JSON{bracket\_certificate}
\]
is present, and the bracket certificate has
\[
        \JSON{status = proved}.
\]
This is the minimal non-debatable audit condition for the certificate.

\section{Repository reproduction commands}

The complete reproduction procedure is as follows.  First, clone the public
repository:
\[
        \code{git clone https://github.com/vicblab/RH-inner-strip-certifier}.
\]
Then enter the repository and create a fresh environment:
\[
        \code{cd RH-inner-strip-certifier},
\]
\[
        \code{python -m venv .venv},
\]
\[
        \code{source .venv/bin/activate},
\]
\[
        \code{python -m pip install --upgrade pip},
\]
\[
        \code{python -m pip install python-flint}.
\]
The full certificate is reproduced by
\[
        \code{./run\_all\_certification.sh}.
\]
The final output file is
\[
        \code{INNER\_STRIP\_CERTIFICATE\_SAFE.json}.
\]
The height threshold is reproduced by
\[
        \code{python tools/height\_threshold.py}.
\]
The repository README gives the same tutorial-level instructions and records
the expected final constants and final JSON status fields. \cite{InnerRepo}

The certification documented in this paper should therefore be read as a
reproducible numerical and interval-arithmetic certificate for the structural
quotient obstruction interface.  The analytic proof of that interface is
external to this certificate paper; the present document verifies the constants,
normalization ledger, bracket lower bound, and height computation used by the
inner-strip obstruction.

\section{Bibliographic references}

The following references are used for the analytic and computational context
of the certification.  The repository reference is the public reproduction
artifact for the certificate.  The remaining references are standard references
for the Riemann zeta function, asymptotic analysis, interval arithmetic, and
Arb-style ball arithmetic.

\begin{thebibliography}{99}

\bibitem{Titchmarsh}
E.~C. Titchmarsh,
\emph{The Theory of the Riemann Zeta-Function},
2nd ed., revised by D.~R. Heath-Brown,
Oxford University Press, 1986.

\bibitem{Edwards}
H.~M. Edwards,
\emph{Riemann's Zeta Function},
Dover Publications, 2001.

\bibitem{Olver}
F.~W.~J. Olver,
\emph{Asymptotics and Special Functions},
AKP Classics, A K Peters, 1997.

\bibitem{WhittakerWatson}
E.~T. Whittaker and G.~N. Watson,
\emph{A Course of Modern Analysis},
4th ed.,
Cambridge University Press, 1927.

\bibitem{Moore}
R.~E. Moore,
\emph{Interval Analysis},
Prentice-Hall, 1966.

\bibitem{Tucker}
W.~Tucker,
\emph{Validated Numerics: A Short Introduction to Rigorous Computations},
Princeton University Press, 2011.

\bibitem{Arb}
F.~Johansson,
\emph{Arb: Efficient Arbitrary-Precision Midpoint-Radius Interval Arithmetic},
IEEE Transactions on Computers 66, no.~8, 2017, pp.~1281--1292.

\bibitem{FLINT}
W.~Hart,
\emph{Fast Library for Number Theory: An Introduction},
Mathematical Software -- ICMS 2010,
Lecture Notes in Computer Science, vol.~6327,
Springer, 2010, pp.~88--91.
\bibitem{OuterCompanion}
V.~Blanco Bataller,
\emph{Certified Outer-Strip Quotient Obstruction for the Riemann Zeta Function},
companion certification manuscript.

\bibitem{OuterRepo}
V.~Blanco Bataller,
\emph{RH-outer-strip-certifier},
certifier repository for the outer-strip residual constants and high-height
outer-strip exclusion.

\bibitem{InnerRepo}
V.~Blanco Bataller,
\emph{RH-inner-strip-certifier},
GitHub repository,
\[
        \url{https://github.com/vicblab/RH-inner-strip-certifier}.
\]
The repository contains the scripts, sector certificates, reproduction
pipeline, and height-threshold utility for the certified inner-strip constants.
The README records the target strip, bracket parameters, expected constants,
one-command reproduction command, and final JSON status requirements.
\end{thebibliography}

\end{document}


